%% 第24回編集委員会（物理第3巻）：組 波W4
%% p3-07.tex  波-7　波が運ぶエネルギー：インピーダンスと反射（★3〜5・3段）
%% 中身は計画書 tools/phys3-plan.md の B「波-7」。公式の持ち主は A（A46・A56）。
%%   分散関係と群速度（A57）は波-5 ⑦で作る（依頼のメッセージ）。この章の⑥は、波-5 ⑦を指して「エネルギーは群速度で進む」ことだけを使う。
%%   固定端・自由端（A16）の持ち主は波-6 ①。この章の⑤は、その一般化（Z₂→∞・Z₂→0）として指す。
%%   球面波の解 Rp＝F(R−ct)（計画書 C-6-1）は⑦で作る（8-12⑦の球のラプラシアンを使う）。波-9 ⑤はこれを指せる。
%% 数値はすべて Python で計算した：checks/物理第3巻/波-7/calc_p3_07.py（出力 calc_p3_07_out.txt）、
%%   sheet_p3_07.py（④の表計算。s＝1 の1歩の式とつなぎ目の式）、checks/物理第3巻/波-8/kanji_p3_07_08.py（常用漢字の確かめ）
%% 記号（この章）：波の速さ c、張力 F_T、線密度 μ、振幅 A、角振動数 ω、振動数 f、波数 k、
%%   インピーダンス Z、媒質の粒の速さ u（速さの振幅 u₀）、音圧 p（振幅 p₀・実効値 p_e）、密度 ρ、仕事率 P、強さ I、
%%   1 m あたりの平均のエネルギー \overline{E}、音源からの距離 R、
%%   振幅（変位・速さ）の反射係数 r・透過係数 τ（時間 t とぶつかるので、透過係数を t と書かない。計画書 C-1 の案）、
%%   吸収の定数 α（1/m）、群速度 c_g（波-5 ⑦）、ばね定数 k_s（波数 k とぶつかるため。波-5 ⑦と同じ）、
%%   ダンパーの係数 b（減衰係数 c は第1部の記号なので、波の速さ c と同じ章では使わない）、屈折率 n（⑨だけ）。
%%   電気の電流は記号で書かない（強さ I とぶつかるため）。
%% 出典（第24回の追記：書くときに確かめた。11月5日。WebSearch で一次資料を引く資料を2つ以上照らし合わせた）
%%   ・空気（20 ℃・1 atm）1.2046 kg/m³・343.3 m/s、水（20 ℃）998.2 kg/m³・1482.3 m/s：CoolProp 8.0（IAPWS-95 と空気の状態式）で計算。
%%       表1 の空気の密度 1.204 は、理想気体の値 101325÷(287.05×293.15)＝1.2041 kg/m³（修3-4：P4-22。ρc は 1.204×343≈413）。
%%       水は NIST Chemistry WebBook（Thermophysical Properties of Fluid Systems、20 ℃・1 atm：998.21 kg/m³・1482.3 m/s）でも確かめた（curl で読んだ）
%%   ・体の軟らかい組織 密度 1060・音速 1540・Z 1.63×10⁶、骨 Z 7.8×10⁶：radiopaedia "Acoustic impedance"、
%%       Stanford RAD225 の講義資料（web.stanford.edu/class/rad225/subdir/Class3.pdf）、cyberphysics の表
%%   ・体の中の吸収 0.5 dB/(cm·MHz)：Stanford RAD225 の表（軟らかい組織 0.5〜2 とする資料もある。代表の値として使う）
%%   ・ゼリー：FDA 510(k) K102637（Dynarex の超音波ゼリー：1530.1 m/s・1.505 MRayl）、A レベル物理の表（1500 m/s・1.5×10⁶）
%%   ・鋼 約5900 m/s・約4.6×10⁷：超音波探傷の音速表（1020 鋼 5890 m/s・46.1×10⁶）など。圧電セラミックス 約3×10⁷：特許文献で 30〜34 MRayl
%%   ・中耳：鼓膜とあぶみ骨の底の面積比 約17〜20、てこ 約1.3（von Békésy の値）、面積比とてこの比は目安で振動数で変わる：
%%       PMC4718164（中耳の働きのレビュー）、米国特許 US4655776 の説明。鼓膜 55 mm²・あぶみ骨 3.2 mm² の数は一次資料で確かめられなかったので、本文では使わない
%%   ・フッ化マグネシウムの屈折率 約1.38、ガラスの反射が4 %→1.2〜1.5 %：findlight.net、syronoptics.com、photonics.com
%%   ・太陽の放射 3.828×10²⁶ W、太陽定数 1361 W/m²：IAU 2015 Resolution B3（arXiv:1510.07674 の本文を pdftotext で読んで確かめた）
%%   ・深い水の波 ω＝√(gk)、群速度＝位相速度の半分：流体力学の教科書の式（Rabaud & Moisy, PRL 110, 214503 (2013)、arXiv:1304.2653 の本文でも確かめた）
%%   ・超音波の検査の歴史：ハウリーらの水につける装置（1950年代、B-29 の銃座を使った水槽）、ドナルドらの1958年 Lancet の論文
%%       （金属の探傷器を借りた）：ob-ultrasound.net、グラスゴー大学のニュース（2008年）
%%   ・ドゥシックの1942年の透過法の実験（修3-4、11月5日に確かめた。手引き 2章の13）：一次資料は K. T. Dussik,
%%       "Über die Möglichkeit, hochfrequente mechanische Schwingungen als diagnostisches Hilfsmittel zu verwerten",
%%       Z. ges. Neurol. Psychiat. 174, 153–168 (1942)。書誌と年を、京都大学の外科の雑誌の論文（1964年。「1942年 Dussik, K. T. の透過法による実験以来」、
%%       repository.kulib.kyoto-u.ac.jp）、医学書院の雑誌の抄録（webview.isho.jp 10.11477/mf.1407204834：1942年に透過法の実験から脳の病気の診断の可能性を指摘）、
%%       Springer の論文の文献表（10.1007/BF02474132）の3つで照らし合わせた（WebSearch。どのサイトも curl では届かなかった）。本文は読んでいない
%% 〔要確認〕（確かめられなかった値）
%%   ・校種：「波のエネルギーは振幅の2乗・振動数の2乗に比例」を高校物理（発展）とした（計画書 A46 の案）。整合の組が、高等学校学習指導要領（平成30年告示）
%%     解説 理科編（文部科学省 https://www.mext.go.jp/content/20250311-mxt_kyoiku02-100002620_05.pdf）で照らした：「物理基礎」(2)(ア)㋑ に
%%     「波がもつエネルギーに触れることも考えられる」とあり、式は挙げていないので（発展）のままにした
\makeatletter
% 超音波のプローブ（横から）  \rxhprobe{(下のまん中)}
\newcommand\rxhprobe[1]{\begin{scope}[shift={#1}]
  \filldraw[fill=bkGray!25, draw=bkInk!70, line width=0.6pt, rounded corners=1.2mm] (-5,0) -- (-6,8) -- (-3.5,18) -- (3.5,18) -- (6,8) -- (5,0) -- cycle;
  \fill[bkNavy!70] (-5,0) rectangle (5,1.2);
  \draw[bkInk!70, line width=1.2pt] (0,18) .. controls (0,22) and (3,24) .. (6,25);\end{scope}}
% 弦の上の三角の山  \rxhpulse{(左はし)}{はば}{高さ}{色}
\newcommand\rxhpulse[4]{\draw[#4, line width=1.3pt] #1 -- ++(#2/2,#3) -- ++(#2/2,-#3);}
% 太い矢印（エネルギーの流れ）  \rxhflow{(始め)}{(終わり)}
\newcommand\rxhflow[2]{\draw[bkOrange!85!black, line width=1.6pt, -{Stealth[length=3.2mm, width=2.6mm]}] #1 -- #2;}
\makeatother

\phchapter{7}{波が運ぶエネルギー：\\インピーダンスと反射}{超音波の検査で、なぜおなかにゼリーをぬる？}{\bklevelrange{3}{5}}{波が運ぶエネルギー：インピーダンスと反射}
\bkmichishirube{3-12,6-9,6-10,8-12,力-4,波-5,波-6}{波-9,波-10,波-13,電-13}

\begin{goal}
\begin{itemize}[label={\color{bkTheme}●}, leftmargin=1.6\zw, itemsep=0pt, topsep=1pt]
\item \bklevel{3}\enspace 波が運ぶエネルギーと強さを、振幅・振動数・距離から求められる。
\item \bklevel{3}\enspace インピーダンスの比から、境目ではね返る割合と通りぬける割合を求められる。
\item \bklevel{3}\enspace 表計算で、2本の弦をつないだときの山のはね返りと通りぬけを追える。
\item \bklevel{5}\enspace 波の式と境界条件から反射係数を導き、固定端・自由端・はね返らない端を1つの式で説明できる。
\item \bklevel{5}\enspace 吸収による弱まり、整合層、光の反射率を、インピーダンスと dB で見積もれる。
\end{itemize}
\end{goal}

\phkoushiki{%
\phk{波のエネルギーは、振幅の2乗と振動数の2乗に比例する}{高校物理（発展）}{ph:波-7:energy}
}

\begin{hatena}
\noindent\begin{minipage}[c]{0.54\linewidth}\raggedright
病院で、おなかの中の赤ちゃんや肝臓のようすを見る\textbf{超音波の検査}では、はじめに、おなかに冷たいゼリーをぬります。その上から、手のひらくらいの器具（プローブ）を当てます。プローブは、1秒に数百万回ふるえる、耳に聞こえない高い音（超音波）を出し、体の中からはね返ってきた音で絵を作ります。

ゼリーをぬらずに、乾いたおなかに当てると、絵はほとんど写りません。薄いゼリーが、なぜそんなに大事なのでしょう。
\end{minipage}\hfill
\begin{minipage}[c]{0.43\linewidth}\centering
\begin{tikzpicture}[x=0.84mm, y=0.84mm]
  % 左：ゼリーなし
  \fill[bkOrange!18] (0,0) rectangle (26,10);
  \draw[bkOrange!70!black, line width=0.7pt] (0,10) -- (26,10);
  \fill[bkBlue!10] (7,10) rectangle (17,12);
  \rxhprobe{(12,12)}
  \foreach \xx in {9,12,15} {\draw[bkRed!80!black, line width=0.7pt, -{Stealth[length=1.5mm]}] (\xx,12) -- (\xx,10.4) -- (\xx+1.2,13.6);}
  \node[font=\sffamily\footnotesize, anchor=north] at (13,-1) {ゼリーなし};
  \node[font=\sffamily\footnotesize, anchor=south west, align=left, inner sep=1pt] at (18.6,12.4) {空気の\\すきま};
  % 右：ゼリーあり
  \begin{scope}[shift={(32,0)}]
  \fill[bkOrange!18] (0,0) rectangle (26,10);
  \draw[bkOrange!70!black, line width=0.7pt] (0,10) -- (26,10);
  \fill[bkGreen!45] (5,10) rectangle (19,12);
  \rxhprobe{(12,12)}
  \foreach \xx in {9,12,15} {\draw[bkRed!80!black, line width=0.7pt, -{Stealth[length=1.5mm]}] (\xx,12) -- (\xx,3);}
  \node[font=\sffamily\footnotesize, anchor=north] at (13,-1) {ゼリーあり};
  \node[font=\sffamily\footnotesize, anchor=south west, inner sep=1pt] at (19.6,12.4) {ゼリー};
  \end{scope}
  \node[font=\sffamily\footnotesize, anchor=west, inner sep=1pt] at (0.5,5) {体};
\end{tikzpicture}
\end{minipage}
\end{hatena}

\noindent 波-5・波-6では、波の形の進み方とはね返り方を見ました。波が進んでも、水や空気やひもなどの\textbf{媒質}そのものは、遠くへは運ばれません（波-5 ①）。それでも、海の大きな波は岸の岩を削り、近くの雷の音は窓ガラスをふるわせます。波は\textbf{エネルギー}を運んでいるのです。この章では、波が運ぶエネルギーの量と、波が別の物へ入るときにはね返る割合・通りぬける割合を調べます。その割合は、2つの物の\textbf{インピーダンス}の比で決まり、超音波の検査からメガネのレンズ、テレビのアンテナ線まで、同じ式で説明できます。

\begin{yattemiyou}[つないだひもと、お風呂の中の音]
\textbf{用意するもの}：太さの違う2本のひも（毛糸と、なわとびのなわなど。それぞれ\qty{2}{m}くらい）、お風呂\par\vspace{2pt}
\noindent\begin{minipage}[c]{0.66\linewidth}
\begin{enumerate}[label={\bkcircled[bkGreen]{\arabic*}}, leftmargin=2.4\zw, itemsep=1pt, topsep=0pt]
\item 細いひもと太いなわを結び、床にまっすぐ置く。太いなわの先は友だちが持つ。
\item 細いひものはしを1回すばやく横にふって、山を送る。つなぎ目ではね返る山は、同じ側か、反対側か。
\item 太いなわのはしから、同じように山を送る。
\item お風呂で耳までもぐり、外の人の声と、お湯の中で打ち合わせたスプーンの音を聞き比べる。
\end{enumerate}
\end{minipage}\hfill
\begin{minipage}[c]{0.32\linewidth}\centering
\begin{tikzpicture}[x=0.85mm, y=0.85mm]
  \draw[bkBlue!70!black, line width=0.6pt] (0,6) -- (22,6);
  \draw[bkOrange!70!black, line width=2.2pt] (22,6) -- (46,6);
  \fill[bkInk] (22,6) circle (1.0);
  \rxhpulse{(4,6)}{8}{6}{bkBlue!70!black}
  \draw[bkBlue!70!black, line width=0.6pt, -{Stealth[length=1.8mm]}] (10,14.5) -- (17,14.5);
  \node[font=\sffamily\footnotesize, anchor=north] at (11,4.2) {細いひも};
  \node[font=\sffamily\footnotesize, anchor=north] at (34,4.2) {太いなわ};
  \node[font=\sffamily\footnotesize, anchor=south] at (27,16) {つなぎ目};
  \draw[bkInk!60, line width=0.4pt] (26,15.8) -- (22.3,7.4);
\end{tikzpicture}
\end{minipage}
\end{yattemiyou}

\noindent やってみようの答え：細いひもから送った山は、つなぎ目で一部がはね返ってきますが、そのとき山は\textbf{反対側}（谷）になっています。残りは、低くゆっくりした山になって、太いなわへ進みます。太いなわから送ると、はね返る山は\textbf{同じ側}のままで、細いひもへは、もとより高い山が速く進みます。お風呂の中では、外の人の声はほとんど聞こえないのに、お湯の中で打ち合わせたスプーンの音は、びっくりするほど大きく聞こえます。空気と水の境目では、音がほとんどはね返ってしまい、入っていけないのです。どちらも、この章の④で数を使って説明します。

\begin{omoidasou}
\bkrecall{波-5}{波の速さ $c$、振動数 $f$、波長 $\lambda$ の間に $c=f\lambda$。正弦波 $y=A\sin(\omega t-kx)$、$\omega=2\pi f$、$k=2\pi/\lambda$。張力 $F_{\mathrm T}$・線密度 $\mu$ の弦の波の速さは $c=\sqrt{F_{\mathrm T}/\mu}$。}
\bkrecall{波-6}{弦の動けない端（固定端）では、山は谷になってはね返る。自由に動ける端（自由端）では、山は山のままはね返る。}
\bkrecall{波-1}{単振動のエネルギーは $\tfrac12\times$（ばね定数）$\times A^2$。ばね定数は（質量）$\times\omega^2$ なので、$\tfrac12m\omega^2A^2$ とも書ける。}
\bkrecall{力-4}{仕事率は 力×速さ（$P=Fv$）。単位は W（＝J/s）。}
\bkrecall{3-12}{デシベル＝10×$\log_{10}$（強さの比）。強さが2倍で約\qty{3}{dB}、10倍で\qty{10}{dB}。}
\bkrecall{6-10}{$e^{j\theta}=\cos\theta+j\sin\theta$。波を $Ae^{j(\omega t-kx)}$ と書き、実部を本物の波と見る（波-5 ⑨）。}
\bkrecall{熱-4}{\qty{20}{\celsius}の空気の音の速さは\qty{343}{m/s}（熱-4 ⑥では $a$ と書いた。この巻では $c$）。}
\end{omoidasou}

\noindent この章の記号：反射係数を $r$、透過係数を $\tau$（タウ）と書きます。透過係数を $t$ と書く本もありますが、この章では時間 $t$ と分けるためです。音源からの距離は $R$ と書きます。

%% ======================== 物理の道具 ========================
\bksection{doguu}{物理の道具}

\phstage{me}{波がエネルギーを運ぶようすを見る}{3}

\bksubsection{1}{波はエネルギーを運ぶ}
\label{ph:波-7:sec1}
長いロープのはしを持って、上下にゆすり続けてみましょう。腕が少し疲れてきます。手は、ロープの張力にさからってロープのはしを上げ下げし、仕事をしています（力-4）。その仕事は、どこへ行ったのでしょう。ロープの遠くの部分も、波が着くと上下に動き出します。動いている物は運動エネルギーをもっています。手がした仕事が、波にのって、ロープの先へ運ばれていったのです。

\begin{bkfigure}
\begin{tikzpicture}[x=1mm, y=1mm]
  \draw[bkGray, line width=0.4pt, dashed] (0,0) -- (124,0);
  \draw[bkBlue, line width=1.4pt, domain=0:120, samples=160, smooth] plot (\x, {8*sin(9*\x)});
  \fill[bkRed] (10,8) circle (1.3);
  \fill[bkRed] (20,0) circle (1.3);
  \fill[bkRed] (30,-8) circle (1.3);
  \draw[bkRed!80!black, line width=1pt, -{Stealth[length=2.2mm]}] (20,1.8) -- (20,10.5);
  \node[font=\sffamily\footnotesize, anchor=south] at (8,10.2) {止まる};
  \node[font=\sffamily\footnotesize, anchor=west] at (21,12) {いちばん速い};
  \node[font=\sffamily\footnotesize, anchor=north] at (30,-10) {止まる};
  \draw[bkInk, line width=0.5pt, {Stealth[length=2pt]}-{Stealth[length=2pt]}] (50,0) -- (50,8);
  \node[font=\sffamily\footnotesize, anchor=west] at (51,4) {$A$};
  \draw[bkInk, line width=0.5pt, {Stealth[length=2pt]}-{Stealth[length=2pt]}] (70,-13) -- (110,-13);
  \draw[bkInk, line width=0.5pt] (70,-11) -- (70,-15) (110,-11) -- (110,-15);
  \node[font=\sffamily\footnotesize, anchor=north] at (90,-13.5) {波長 $\lambda$};
  \rxhflow{(66,15)}{(96,15)}
  \node[font=\sffamily\footnotesize, anchor=south] at (81,16.5) {波は速さ $c$ で右へ。エネルギーも右へ};
\end{tikzpicture}
\bkfigcaption{ロープを伝わる正弦波。ロープの1つ1つの部分は、その場で上下に単振動するだけ（山と谷で止まり、まん中でいちばん速い）。動きのエネルギーが、波といっしょに右へ運ばれる。}\label{ph:波-7:fig-rope}
\end{bkfigure}

\phboxhead{ロープ1 m がもつエネルギー}
振幅 $A$、角振動数 $\omega$ の正弦波が、線密度 $\mu$〔kg/m〕のロープを伝わっているとします。ロープを短い長さ $\Delta x$ に区切ると、1つの部分の質量は $\mu\,\Delta x$ です。波が通るあいだ、この部分は、その場で上下に振幅 $A$、角振動数 $\omega$ の単振動をします（波-5 ③）。単振動のエネルギーは $\tfrac12m\omega^2A^2$ でした（波-1 ⑥）。$m$ に $\mu\,\Delta x$ を入れて、$\Delta x$ でわると、ロープ1 m あたりのエネルギー $\overline E$〔J/m〕は
\pheq{\overline E=\frac12\mu\omega^2A^2=2\pi^2\mu f^2A^2}{J/m ＝ $\dfrac{\text{kg}}{\text{m}}$×$\dfrac{1}{\text{s}^2}$×m$^2$}
です（$\omega=2\pi f$ を入れると、$\tfrac12\times(2\pi)^2=2\pi^2$）。\label{ph:波-7:energy}エネルギーは、\textbf{振幅の2乗}と\textbf{振動数の2乗}に比例します。$E$ の上の線は「平均」の印です（わけは、このあとの「せいかくに言うと」）。

\begin{phyomi}[$\overline E=\tfrac12\mu\omega^2A^2$]
「波がもつエネルギーは、媒質の重さ（$\mu$）、ゆれの速さ（$\omega$）、ゆれの大きさ（$A$）の3つで決まる。$\omega$ と $A$ は2乗で効く」。ロープの1つの部分がいちばん速く動くときの速さは $A\omega$ です（波-5 ③）。$\tfrac12\mu(A\omega)^2$ は、「1 m のロープが、いちばん速い速さで動いているときの運動エネルギー」と読めます。
\end{phyomi}

\begin{phbaisuru}{振幅 $A$}
エネルギーは振幅の2乗に比例するので、振幅を2倍にすると4倍、3倍にすると9倍になります。振動数を2倍にしても、同じように4倍です。下のグラフは、例題3のロープ（$\mu=\qty{0.20}{kg/m}$、$f=\qty{2}{Hz}$）の1 m あたりのエネルギーです。
\par\smallskip\noindent\begin{tikzpicture}
\begin{axis}[bkplot, width=96mm, height=42mm, xmin=0, xmax=0.115, ymin=0, ymax=0.19, xtick={0.05,0.10}, xticklabel style={/pgf/number format/fixed}, ytick={0.0395,0.158}, yticklabels={0.039,0.16}, xlabel={振幅 $A$（m）}, ylabel={$\overline E$（J/m）}]
  \addplot[bkcurve, domain=0:0.11, samples=60] {0.5*0.2*(4*pi)^2*x^2};
  \fill[bkRed] (axis cs:0.05,0.0395) circle (1.8pt) (axis cs:0.10,0.158) circle (1.8pt);
  \node[font=\sffamily\footnotesize, anchor=north west, text=bkRed!75!black] at (axis cs:0.064,0.07) {振幅2倍で4倍};
\end{axis}
\end{tikzpicture}
\end{phbaisuru}

\begin{phmijika}[海の波と台風の高波]
海の波の高さ（山から谷まで。振幅の2倍）が\qty{1}{m}の日と\qty{2}{m}の日を比べると、高さが2倍なので、波がもつエネルギーは4倍です。台風が近づいて高さ\qty{6}{m}の波が来ると、\qty{1}{m}の波の36倍のエネルギーが、海岸の堤防にぶつかります。高い波が、見た目の高さ以上に危ないのは、このためです。
\end{phmijika}

\begin{seikaku}
単振動するおもりでは、運動エネルギーと位置エネルギーが入れかわり、和はいつも $\tfrac12m\omega^2A^2$ です。進む波のロープの1つの部分は、少し違います。図\ref{ph:波-7:fig-rope}で、まん中を通る部分は、いちばん速いうえに、ロープがいちばん斜めになって伸ばされています。山や谷の部分は、止まっていて、ロープも伸びていません。進む波では、運動エネルギーと伸びのエネルギーが\textbf{同時に}大きくなり、同時に0になるのです。それでも、1 m のロープについて平均すると、エネルギーはちょうど $\tfrac12\mu\omega^2A^2$ になります（伸びのエネルギーを式で書いて確かめる計算は、専門の本で）。
\end{seikaku}

\phboxhead{エネルギーが流れる速さ：仕事率}
波は速さ $c$ で進みます。1秒たつと、波は $c$〔m〕先へ進むので、1秒に、ある場所を通りすぎるエネルギーは、長さ $c$ のロープの分です。これが、波が運ぶ\textbf{仕事率} $P$〔W〕です。
\pheq{P=\overline E\,c=\frac12\mu c\,\omega^2A^2}{W ＝ $\dfrac{\text{J}}{\text{m}}$×$\dfrac{\text{m}}{\text{s}}$}
\label{ph:波-7:power}ロープをゆする手は、1秒に $P$ だけの仕事をし続けなければなりません（例題3）。

\bksubsection{2}{強さ：広がると弱くなる}
\label{ph:波-7:sec2}
音や光は、ロープと違って、空間に広がって進みます。そこで、波の進む向きに直角な面 1 m$^2$ を、1秒に通るエネルギーを考えます。これを波の\textbf{強さ}といい、$I$〔W/m$^2$〕と書きます。\bkidx{なみのつよさ}{強さ（波）}

\begin{bkfigure}
\begin{tikzpicture}[x=1mm, y=1mm]
  \fill[bkYellow!70!black] (0,0) circle (1.6);
  \draw[bkBlue!70!black, line width=0.8pt] (0,0) circle (12);
  \draw[bkBlue!70!black, line width=0.8pt] (0,0) circle (24);
  \foreach \a in {20,65,110,155,200,245,290,335} {\draw[bkOrange!85!black, line width=0.9pt, -{Stealth[length=2mm]}] (\a:3) -- (\a:9.5);}
  \fill[bkRed!45] (-4,11.2) rectangle (4,12.8);
  \fill[bkRed!45] (-8,23.2) rectangle (8,24.8);
  \draw[bkInk, line width=0.5pt] (0,0) -- (-35:12);
  \node[font=\sffamily\footnotesize, anchor=north, fill=white, inner sep=1pt] at (-35:15) {$R$};
  \draw[bkInk, line width=0.5pt] (0,0) -- (-140:24);
  \node[font=\sffamily\footnotesize, anchor=north east, fill=white, inner sep=1pt] at (-140:24.5) {$2R$};
  \node[anchor=west, font=\small, text width=74mm, align=left] at (32,4) {%
  音源が1秒に出すエネルギー $P$ は、半径 $R$ の球面の全部（面積 $4\pi R^2$）を通る。\par\vspace{2pt}
  半径が2倍の球面は、面積が4倍。同じ $P$ が4倍の面積に分かれるので、1 m$^2$ あたり（強さ）は $\tfrac14$。\par\vspace{2pt}
  色をつけた帯は、同じ向きに進むエネルギーが通る面。遠くでは広がって、面積が4倍になる。};
\end{tikzpicture}
\bkfigcaption{点の音源から広がる波。どの球面を通るエネルギーも、1秒あたり $P$ で同じ（とちゅうで吸収されなければ）。}\label{ph:波-7:fig-kyuu}
\end{bkfigure}

\noindent 小さな音源が、1秒に $P$〔W〕のエネルギーを、すべての向きに同じように出しているとします。とちゅうで吸収されなければ、そのエネルギーは、音源を囲む半径 $R$ の球面を全部通りぬけます。球面の面積は $4\pi R^2$ なので、
\pheq{I=\frac{P}{4\pi R^2}}{W/m$^2$ ＝ W÷m$^2$}
です。\label{ph:波-7:gyaku2}強さは\textbf{距離の2乗に反比例}します（\textbf{逆2乗の法則}）。距離が2倍になると強さは $\tfrac14$、10倍になると $\tfrac{1}{100}$ です。デシベルで書くと、$10\log_{10}4\approx6.0$ なので、距離が2倍になるごとに約\qty{6}{dB}下がります（3-12⑦）。\bkidx{ぎゃく2じょうのほうそく}{逆2乗の法則}

\begin{phbaisuru}{距離 $R$}
強さは $\tfrac14$ になり、dB で表した音の大きさは約\qty{6}{dB}下がります。下のグラフは、1秒に\qty{0.010}{W}の音を出す小さなスピーカー（例題10）からの強さです。
\par\smallskip\noindent\begin{tikzpicture}
\begin{axis}[bkplot, width=96mm, height=42mm, xmin=0, xmax=4.6, ymin=0, ymax=1.05, xtick={1,2,3,4}, ytick={0.199,0.796}, yticklabels={0.20,0.80}, xlabel={距離 $R$（m）}, ylabel={$I$（mW/m$^2$）}]
  \addplot[bkcurve, domain=0.88:4.4, samples=80] {10/(4*pi*x^2)};
  \fill[bkRed] (axis cs:1,0.796) circle (1.8pt) (axis cs:2,0.199) circle (1.8pt) (axis cs:4,0.0497) circle (1.8pt);
  \node[font=\sffamily\footnotesize, anchor=south west, text=bkRed!75!black] at (axis cs:2.1,0.26) {距離2倍で $\tfrac14$（$-6$ dB）};
\end{axis}
\end{tikzpicture}
\end{phbaisuru}

\begin{phmijika}[太陽の光の強さ]
太陽は、1秒に約 $3.83\times10^{26}$\,J の光を、四方八方に出しています。地球は太陽から約 $1.496\times10^{11}$\,m 離れているので、地球の所での強さは
\[I=\frac{3.83\times10^{26}}{4\pi\times(1.496\times10^{11})^2}\approx\qty{1360}{W/m^2}\]
です。人工衛星ではかった値（約\qty{1361}{W/m^2}）とよく合います。太陽から地球の約1.5倍離れた火星では、強さは約 $1\div1.5^2\approx0.44$ 倍しかありません。
\end{phmijika}

\phboxhead{広がり方で、弱まり方が変わる}
車がたくさん走る道路や、長い列車は、点ではなく\textbf{線}の音源です。音は円柱の形に広がり、円柱の側面の面積は距離 $R$ に比例するので、強さは $1/R$ に比例します。距離が2倍でも、約\qty{3}{dB}しか下がりません。管の中を進む音（聴診器、昔の船の伝声管）や、糸電話の糸の振動は、広がらないので、ほとんど弱まらずに遠くまで届きます。

\begin{seikaku}
逆2乗の法則は、音源が距離に比べて小さく、周りに音をはね返す物がないときの式です。音源のすぐ近くや、部屋の中（壁ではね返った音が重なる）では、そのとおりには下がりません（波-9）。
\end{seikaku}

\Needspace{9\baselineskip}\phstage{te}{インピーダンスで、はね返りを求める}{3}

\bksubsection{3}{インピーダンス：押す力÷動く速さ}
\label{ph:波-7:sec3}
とても長いロープ（はね返りが戻ってこないくらい長い）のはしを持ち、一定の速さ $u$ で上へ動かし始めたとします。はしが上がると、ロープは折れ曲がり、その折れ目が速さ $c$ で右へ進んでいきます（波-5）。

\begin{bkfigure}
\begin{tikzpicture}[x=1mm, y=1mm]
  \draw[bkGray, line width=0.4pt, dashed] (0,0) -- (76,0);
  \draw[bkBlue, line width=1.4pt] (0,10) -- (42,0) -- (76,0);
  \fill[bkInk] (0,10) circle (1.1);
  \draw[bkRed!80!black, line width=1pt, -{Stealth[length=2.2mm]}] (-3,3) -- (-3,12);
  \node[font=\sffamily\footnotesize, anchor=south] at (-3,12.6) {速さ $u$};
  \draw[bkInk, line width=0.5pt, {Stealth[length=2pt]}-{Stealth[length=2pt]}] (2.5,0) -- (2.5,10);
  \node[font=\sffamily\footnotesize, anchor=west, fill=white, inner sep=1pt] at (3.5,3.2) {$u\,\Delta t$};
  \draw[bkInk, line width=0.5pt, {Stealth[length=2pt]}-{Stealth[length=2pt]}] (0,-4) -- (42,-4);
  \node[font=\sffamily\footnotesize, anchor=north] at (21,-4.6) {$c\,\Delta t$（折れ目が進んだ距離）};
  \draw[bkOrange!85!black, line width=1pt, -{Stealth[length=2.2mm]}] (42,0) -- (52,0);
  \node[font=\sffamily\footnotesize, anchor=south west] at (43,1) {速さ $c$};
  \draw[bkGreen!60!black, line width=1.1pt, -{Stealth[length=2.4mm]}] (0,10) -- (15,6.43);
  \node[font=\sffamily\footnotesize, anchor=south west, text=bkGreen!50!black] at (16,7.2) {張力 $F_{\mathrm T}$};
  \node[anchor=west, font=\small, text width=44mm, align=left] at (78,6) {%
  斜めの部分の傾きは\par $\dfrac{u\,\Delta t}{c\,\Delta t}=\dfrac{u}{c}$\par\vspace{2pt}
  手は、張力のたての成分 $F_{\mathrm T}\times\dfrac{u}{c}$ を\par 支え続ける。};
\end{tikzpicture}
\bkfigcaption{長いロープのはしを、一定の速さ $u$ で上げていく。時間 $\Delta t$ のあいだに、はしは $u\,\Delta t$ 上がり、折れ目は $c\,\Delta t$ 進む。}\label{ph:波-7:fig-oshi}
\end{bkfigure}

\noindent 時間 $\Delta t$ たつと、はしは $u\,\Delta t$ だけ上がり、折れ目は $c\,\Delta t$ だけ進みます。斜めになった部分の傾きは $u/c$ です。傾きが小さいとき、ロープがはしを引く張力 $F_{\mathrm T}$ のたての成分は、$F_{\mathrm T}\times$（傾き）です。手は、このたての力を上へ押し返しながら、はしを動かし続けなければなりません。手の力 $F$ は、
\[F=F_{\mathrm T}\times\frac{u}{c}=\frac{F_{\mathrm T}}{c}\,u\]
で、\textbf{速さ $u$ に比例}します。比例の定数は、$c=\sqrt{F_{\mathrm T}/\mu}$ を使うと、
\[\frac{F_{\mathrm T}}{c}=\frac{F_{\mathrm T}}{\sqrt{F_{\mathrm T}/\mu}}=\sqrt{F_{\mathrm T}\mu}=\mu c\]
です。この定数を、ロープの\textbf{インピーダンス}といい、$Z$ と書きます。\bkidx{いんぴーだんす}{インピーダンス}
\pheq{Z=\frac{F}{u}=\sqrt{F_{\mathrm T}\mu}=\mu c}{N·s/m ＝ N÷(m/s)}
\label{ph:波-7:Z}

\phboxhead{もう1つの見方：運動量で数える}
同じ式を、運動量（力-3）でも作れます。時間 $\Delta t$ のあいだに、新しく動き出したロープは長さ $c\,\Delta t$、質量 $\mu c\,\Delta t$ で、速さ0から $u$ になりました。運動量の変化 $\mu c\,\Delta t\times u$ は、手が加えた力積 $F\,\Delta t$ に等しいので、$F=\mu c\,u$。やはり $Z=\mu c$ です。

\begin{phyomi}[$Z=F/u$]
「媒質を速さ $u$ で動かすのに要る力は $Zu$。$Z$ は、その媒質の\textbf{手ごたえ}」。$Z$ が大きい媒質は、同じ速さで動かすのに大きな力が要ります。重い（$\mu$ が大きい）ほど、強く張ってある（$F_{\mathrm T}$ が大きい）ほど、$Z$ は大きくなります。インピーダンス（impedance）は、英語で「じゃまをする」という意味のことばから来ています。
\end{phyomi}

\begin{phbaisuru}[4倍]{張力 $F_{\mathrm T}$}
$Z=\sqrt{F_{\mathrm T}\mu}$ なので、張力を4倍にすると、$Z$ は $\sqrt4=2$ 倍になります。線密度を4倍にしても2倍です。ただし、波の速さ $c=\sqrt{F_{\mathrm T}/\mu}$ は、張力4倍では2倍、線密度4倍では半分です。$Z$ と $c$ は別の量です。
\end{phbaisuru}

\phboxhead{インピーダンスと仕事率}
手が1秒にする仕事は、力×速さ $=Fu=Zu^2$ です（力-4）。正弦波では、はしの速さは $u=A\omega\cos\omega t$ と変わるので、$u^2$ の平均は $\tfrac12A^2\omega^2$ です（$\cos^2$ の平均は $\tfrac12$：4-9⑧）。だから、手がする仕事率の平均は
\pheq{P=\frac12Z\omega^2A^2}{W ＝ $\dfrac{\text{N·s}}{\text{m}}$×$\dfrac{1}{\text{s}^2}$×m$^2$}
です。$Z=\mu c$ を入れると①の $P=\tfrac12\mu c\,\omega^2A^2$ と同じになり、2つの考えが合っています。

\phboxhead{音のインピーダンス}
長い管の中の空気を、ピストンで速さ $u$ で押すときも同じです。時間 $\Delta t$ に、押し縮められた空気の部分は長さ $c\,\Delta t$ だけ伸び、その質量 $\rho Sc\,\Delta t$（$S$ は管の断面積）が速さ $u$ で動き出します。ピストンが押す圧力の増えた分を $p$ とすると、力積 $pS\,\Delta t$ が運動量 $\rho Sc\,\Delta t\times u$ に等しいので、
\pheq{p=\rho c\,u\qquad Z=\frac{p}{u}=\rho c}{Pa·s/m ＝ $\dfrac{\text{kg}}{\text{m}^3}$×$\dfrac{\text{m}}{\text{s}}$}
\label{ph:波-7:onkyo}です。$p$ は大気圧からの圧力のずれで、\textbf{音圧}といいます。\bkidx{おんあつ}{音圧}$Z=\rho c$〔Pa·s/m〕を、その物の\textbf{音響インピーダンス}といいます。\bkidx{おんきょういんぴーだんす}{音響インピーダンス}20 ℃の空気では、$1.204\times343\approx413$ です。

音の強さ（1 m$^2$ あたりの仕事率）は、圧力×速さ $=pu=p^2/(\rho c)$ の平均です。音圧の振幅を $p_0$ とすると、$p^2$ の平均は $\tfrac12p_0^2$ です。$p_0/\sqrt2$ を音圧の\textbf{実効値} $p_{\mathrm e}$ といいます（交流の電圧の実効値と同じ考え）。
\pheq{I=\frac{p_{\mathrm e}^{\,2}}{\rho c}=\frac{p_{\mathrm e}^{\,2}}{Z}}{W/m$^2$ ＝ Pa$^2$÷(Pa·s/m)}
\label{ph:波-7:onI}同じ音圧なら、$Z$ の小さい空気のほうが、水よりずっと大きなエネルギーを運びます（例題2）。

\begin{bkcenter}
\begin{bktable}{colspec={Q[l,m,wd=40mm] Q[c,m,wd=26mm] Q[c,m,wd=24mm] Q[c,m,wd=30mm]}, rowsep=1.5pt, row{1}={bg=bkTheme!10, font=\sffamily\small}}
物 & 密度 $\rho$〔kg/m$^3$〕 & 音速 $c$〔m/s〕 & $Z=\rho c$〔Pa·s/m〕 \\
空気（\qty{20}{\celsius}） & 1.204 & 343 & 413 \\
水（\qty{20}{\celsius}） & 998 & 1482 & $1.48\times10^6$ \\
体の軟らかい組織 & 約1060 & 約1540 & $1.63\times10^6$ \\
超音波検査のゼリー & ― & ― & 約 $1.5\times10^6$ \\
骨 & ― & ― & 約 $7.8\times10^6$ \\
圧電セラミックス & ― & ― & 約 $3\times10^7$ \\
鋼 & 7850 & 約5900 & 約 $4.6\times10^7$ \\
\end{bktable}
\bktabcaption{いろいろな物の音響インピーダンス（代表の値）。空気と体では約4000倍違う。圧電セラミックスは、プローブの中で電気を音に変える部品の材料（波-10）。}\label{ph:波-7:tab-Z}
\end{bkcenter}

\noindent ロープの $Z$ は「力÷速さ」、音の $Z$ は「圧力÷速さ」でした。どちらも「\textbf{押すもの}÷\textbf{流れるもの}」です。電気の線（テレビのアンテナ線などの同軸ケーブル）を伝わる波でも、電圧÷電流が決まった値になり、これを\textbf{特性インピーダンス}といいます（50 Ω や 75 Ω。物理シリーズの『電磁気学』（電-13））。

\bksubsection{4}{はね返りと通りぬけ：2本の弦をつなぐ}
\label{ph:波-7:sec4}
インピーダンス $Z_1$ の弦に、インピーダンス $Z_2$ の弦をつなぎ、左から波を送ります。つなぎ目に着いた波は、一部がはね返り（\textbf{反射}）、残りが右の弦へ進みます（\textbf{透過}）。\bkidx{はんしゃ}{反射}\bkidx{とうか}{透過}その割合を、2つの決まりから求めましょう。

\begin{bkfigure}
\begin{tikzpicture}[x=1mm, y=1mm]
  \draw[bkBlue!70!black, line width=0.7pt] (0,0) -- (60,0);
  \draw[bkOrange!70!black, line width=2.0pt] (60,0) -- (120,0);
  \fill[bkInk] (60,0) circle (1.2);
  \rxhpulse{(14,0)}{14}{10}{bkBlue!80!black}
  \draw[bkBlue!80!black, line width=0.8pt, -{Stealth[length=2mm]}] (16,13) -- (28,13);
  \node[font=\sffamily\footnotesize, anchor=south] at (22,14) {来た波（速さ $u_{\text{来}}$）};
  \rxhpulse{(36,0)}{14}{-4}{bkRed!80!black}
  \draw[bkRed!80!black, line width=0.8pt, -{Stealth[length=2mm]}] (52,-7) -- (40,-7);
  \node[font=\sffamily\footnotesize, anchor=north] at (44,-8) {はね返った波 $u_{\text{返}}$};
  \rxhpulse{(74,0)}{8}{6.5}{bkOrange!80!black}
  \draw[bkOrange!80!black, line width=0.8pt, -{Stealth[length=2mm]}] (74,10) -- (84,10);
  \node[font=\sffamily\footnotesize, anchor=south west] at (73,11) {通りぬけた波 $u_{\text{通}}$};
  \node[font=\sffamily\footnotesize, anchor=north] at (24,-2) {弦1（$Z_1$）};
  \node[font=\sffamily\footnotesize, anchor=north] at (102,-2) {弦2（$Z_2$）};
  \node[font=\sffamily\footnotesize, anchor=south] at (60,2) {つなぎ目};
\end{tikzpicture}
\bkfigcaption{2本の弦のつなぎ目での反射と透過（$Z_2>Z_1$ のとき）。$u$ は、それぞれの波がつなぎ目を動かす速さ（上向きを正）。}\label{ph:波-7:fig-tsunagi}
\end{bkfigure}

\phboxhead{決まり1：つなぎ目は1つの点}
つなぎ目は、左の弦のはしでもあり、右の弦のはしでもあります。だから、つなぎ目の速さは、左から見ても右から見ても同じです。左の弦では、来た波とはね返った波が重なっています（波-5 ⑥）。
\[u_{\text{来}}+u_{\text{返}}=u_{\text{通}}\]

\phboxhead{決まり2：つなぎ目の力はつり合う}
結び目はとても軽いので、左の弦がつなぎ目を引く力と、右の弦が引く力は、つり合っていなければなりません（質量0の物に力が残ると、加速度が無限になってしまう）。③で、右へ進む波は、速さ $u$ のとき大きさ $Zu$ の力を伝えることがわかりました。左へ進む波では、ロープの傾きが逆になるので、力の向きも逆の $-Zu$ です。左の弦の力（来た波とはね返った波の分）と、右の弦の力が等しいので、
\[Z_1u_{\text{来}}-Z_1u_{\text{返}}=Z_2u_{\text{通}}\]

\phboxhead{2つの式を解く}
決まり1の $u_{\text{通}}$ を決まり2に入れると、$Z_1u_{\text{来}}-Z_1u_{\text{返}}=Z_2u_{\text{来}}+Z_2u_{\text{返}}$。$u_{\text{返}}$ を左に集めると、$(Z_1+Z_2)u_{\text{返}}=(Z_1-Z_2)u_{\text{来}}$ です。だから、
\pheq{r=\frac{u_{\text{返}}}{u_{\text{来}}}=\frac{Z_1-Z_2}{Z_1+Z_2}\qquad\tau=\frac{u_{\text{通}}}{u_{\text{来}}}=\frac{2Z_1}{Z_1+Z_2}}{単位なし}
\label{ph:波-7:rt}$r$ を\textbf{反射係数}、$\tau$ を\textbf{透過係数}といいます（時刻の $t$ と分けて $\tau$ と書きます。波-4 の時定数 $\tau$ とは別の量です）。\bkidx{はんしゃけいすう}{反射係数}\bkidx{とうかけいすう}{透過係数}ゆれの速さの比は、ゆれの高さ（変位）の比と同じなので、$r$ と $\tau$ は、山の高さの比でもあります。$1+r=\tau$ も成り立っています（決まり1）。

\phboxhead{2つの数え方：速さで数えるか、押すもので数えるか}
\label{ph:波-7:fugou}この章の $r$ と $\tau$ は、ゆれの\textbf{速さ（変位）}で数えた値です。反射係数は、「押すもの」（弦では力、音では音圧 $p$、電気では電圧）で数えることもできます。③で見たように、右へ進む波の力は $+Zu$、左へ進む波の力は $-Zu$ でした。はね返った波は左へ進むので、それが伝える力は $-Z_1u_{\text{返}}=-Z_1\,r\,u_{\text{来}}=-r\times(Z_1u_{\text{来}})$ です。来た波の力 $Z_1u_{\text{来}}$ との比は $-r$ になります。
\pheq{r_{\text{押}}=\frac{\text{はね返った波の力（音圧）}}{\text{来た波の力（音圧）}}=-r=\frac{Z_2-Z_1}{Z_1+Z_2}}{単位なし}
進む向きが逆になると、同じ速さでも押す向きが逆になる（速さと力の比 $Z$ の符号が変わる）ので、\textbf{2つの数え方では、反射係数の符号が逆}になるのです。線密度4倍の弦（$Z_2=2Z_1$）では、速さ・変位で数えると $r=-\tfrac13$（山が谷になってはね返る）、力で数えると $r_{\text{押}}=+\tfrac13$ です。音響や超音波の本（波-10）は、ふつう音圧で数えた $r_{\text{押}}=(Z_2-Z_1)/(Z_2+Z_1)$ を使います。どちらで数えても、はね返るエネルギーの割合は $r^2=r_{\text{押}}^{\,2}$ で同じです。

\phboxhead{エネルギーで見る：はね返る割合と通りぬける割合}
仕事率は $\tfrac12Z\omega^2A^2$ でした（③）。来た波・はね返った波・通りぬけた波は、どれも同じ振動数です（つなぎ目は1つの点で、両側を同じリズムでゆらすから）。だから、エネルギーの割合は、
\pheq{\text{はね返る割合}=r^2\qquad\text{通りぬける割合}=\frac{Z_2}{Z_1}\tau^2=\frac{4Z_1Z_2}{(Z_1+Z_2)^2}}{単位なし}
\label{ph:波-7:wariai}です。2つをたすと、ちょうど1になります（エネルギーの保存。例題4で確かめる）。

\begin{phyomi}[$r=\dfrac{Z_1-Z_2}{Z_1+Z_2}$]
「はね返りは、インピーダンスの\textbf{違い}で決まる。同じなら、はね返らない」。
\begin{itemize}[leftmargin=1.2\zw, itemsep=1pt, topsep=1pt, label={\color{bkTheme}●}]
\item $Z_2>Z_1$（重い・かたい物へ）：$r<0$。はね返る山は\textbf{反対側}になる。
\item $Z_2<Z_1$（軽い・やわらかい物へ）：$r>0$。はね返る山は\textbf{同じ側}のまま。
\item $Z_2=Z_1$：$r=0$。つなぎ目があっても、まったくはね返らない。
\end{itemize}
やってみようの細いひも（$Z$ が小さい）→太いなわ（$Z$ が大きい）は1つめ、太いなわ→細いひもは2つめでした。
\end{phyomi}

\begin{bkfigure}
\begin{tikzpicture}
\begin{axis}[bkplot, axis lines=left, width=118mm, height=56mm, xmode=log, xmin=0.01, xmax=100, ymin=-1.1, ymax=1.1, xtick={0.01,0.1,1,10,100}, xticklabels={0.01,0.1,1,10,100}, ytick={-1,-0.5,0,0.5,1}, xlabel={$Z_2/Z_1$}, ylabel={}, legend style={draw=none, fill=white, font=\sffamily\footnotesize, at={(0.02,0.04)}, anchor=south west}]
  \addplot[bkcurve, domain=0.01:100, samples=120] {(1-x)/(1+x)};
  \addlegendentry{反射係数 $r$}
  \addplot[bkOrange, line width=1.5pt, dash pattern=on 4pt off 2pt, domain=0.01:100, samples=120] {((1-x)/(1+x))^2};
  \addlegendentry{はね返る割合 $r^2$}
  \fill[bkRed] (axis cs:2,-0.3333) circle (1.8pt) (axis cs:2,0.1111) circle (1.8pt);
  \node[font=\sffamily\footnotesize, anchor=south west, text=bkRed!75!black, fill=white, inner sep=1pt] at (axis cs:2.3,0.15) {$Z_2=2Z_1$：$r=-\tfrac13$、$r^2=\tfrac19$};
\end{axis}
\end{tikzpicture}
\bkfigcaption{インピーダンスの比と、反射係数 $r$（実線）・はね返る割合 $r^2$（破線）。横軸は10倍ごとに同じ長さ（3-12 の対数の目もり）。比が1から10倍以上ずれると、半分以上のエネルギーがはね返る。}\label{ph:波-7:fig-rgraph}
\end{bkfigure}

\phboxhead{表計算で確かめる}
2本の弦を、表計算で動かしてみましょう。左の弦は、張力\qty{80}{N}・線密度\qty{0.20}{kg/m}（$c_1=\qty{20}{m/s}$、$Z_1=\qty{4.0}{N.s/m}$）、右の弦は、張力が同じで線密度が4倍の\qty{0.80}{kg/m}（$c_2=\qty{10}{m/s}$、$Z_2=\qty{8.0}{N.s/m}$）です。左の弦を\qty{0.1}{m}ごと、右の弦を\qty{0.05}{m}ごとの点に分け、$\Delta t=\qty{0.005}{s}$ ずつ時間を進めます。どちらの弦でも、1歩で波がちょうど1つとなりの点へ進む（$s=c\,\Delta t/\Delta x=1$）ように選んであります。

表計算の1行が1つの時刻、1つの列が弦の1つの点です。波-5 ⑧の1歩の式は、$s=1$ のとき、とても簡単になります。
\begin{bkcenter}
次の時刻の値＝（今の時刻の、左となり＋右となり）－（1つ前の時刻の、同じ点）
\end{bkcenter}
つなぎ目の点だけは、両側の弦の重さと引っぱりを入れた、次の式にします（わけは、このあとの「せいかくに言うと」）。
\begin{bkcenter}
つなぎ目の次の値＝$-$（1つ前の時刻の値）$+\dfrac{2Z_2}{Z_1+Z_2}$×（今の右となり）$+\dfrac{2Z_1}{Z_1+Z_2}$×（今の左となり）
\end{bkcenter}
$Z_2=2Z_1$ なので、右となりに $\tfrac43$、左となりに $\tfrac23$ をかけます。はじめの2行に、高さ\qty{3}{cm}の三角の山を、1マス右へずらして入れておくと、山は右へ進み始めます。

\begin{bkcenter}
\begin{bktable}{colspec={Q[c,m,wd=17mm] *{6}{Q[c,m,wd=12.5mm]}}, rowsep=1pt, row{1,2}={bg=bkTheme!10, font=\sffamily\footnotesize}, cells={font=\sffamily\footnotesize}, column{4}={bg=bkYellow!20}, cell{1}{2}={c=6}{c}, cell{1}{1}={r=2}{c}}
$t$〔s〕 & 位置 $x$〔m〕 & & & & & \\
 & 1.8 & 1.9 & 2.0 & 2.05 & 2.10 & 2.15 \\
0.065 & 3.00 & 2.00 & 0.67 & 0 & 0 & 0 \\
0.070 & 2.00 & 2.67 & 1.33 & 0.67 & 0 & 0 \\
0.075 & 0.67 & 1.33 & 2.00 & 1.33 & 0.67 & 0 \\
0.080 & $-0.67$ & 0 & 1.33 & 2.00 & 1.33 & 0.67 \\
0.085 & $-1.00$ & $-0.67$ & 0.67 & 1.33 & 2.00 & 1.33 \\
0.090 & $-0.67$ & $-0.33$ & 0 & 0.67 & 1.33 & 2.00 \\
\end{bktable}
\bktabcaption{表計算の一部（つなぎ目の近く。単位は cm）。色をつけた列 $x=\qty{2.0}{m}$ がつなぎ目。たとえば $t=0.075$ の行の $x=2.0$ は、$-0.67+\tfrac43\times0.67+\tfrac23\times2.67\approx2.00$。}\label{ph:波-7:tab-sheet}
\end{bkcenter}

\begin{bkfigure}
\begin{tikzpicture}
\begin{axis}[bkplot, width=124mm, height=50mm, xmin=0, xmax=3.1, ymin=-1.6, ymax=3.6, xtick={0.5,1,1.5,2,2.5,3}, ytick={-1,1,2,3}, xlabel={$x$（m）}, ylabel={$y$（cm）}, xticklabel style={/pgf/number format/fixed}]
  \fill[bkOrange!12] (axis cs:2,-1.6) rectangle (axis cs:3.1,3.6);
  \draw[bkGray, line width=0.9pt, line join=round] (axis cs:0,0) -- (axis cs:1.0,0) -- (axis cs:1.1,1) -- (axis cs:1.2,2) -- (axis cs:1.3,3) -- (axis cs:1.4,2) -- (axis cs:1.5,1) -- (axis cs:1.6,0) -- (axis cs:3.0,0);
  \draw[bkNavy, line width=1.2pt, dash pattern=on 4pt off 2pt, line join=round] (axis cs:0,0) -- (axis cs:1.6,0) -- (axis cs:1.7,-0.33) -- (axis cs:1.8,-0.67) -- (axis cs:1.9,0) -- (axis cs:2.0,1.33) -- (axis cs:2.05,2.0) -- (axis cs:2.10,1.33) -- (axis cs:2.15,0.67) -- (axis cs:2.2,0) -- (axis cs:3.0,0);
  \draw[bkRed!80!black, line width=1.4pt, line join=round] (axis cs:0,0) -- (axis cs:0.8,0) -- (axis cs:0.9,-0.33) -- (axis cs:1.0,-0.67) -- (axis cs:1.1,-1.0) -- (axis cs:1.2,-0.67) -- (axis cs:1.3,-0.33) -- (axis cs:1.4,0) -- (axis cs:2.3,0) -- (axis cs:2.35,0.67) -- (axis cs:2.4,1.33) -- (axis cs:2.45,2.0) -- (axis cs:2.5,1.33) -- (axis cs:2.55,0.67) -- (axis cs:2.6,0) -- (axis cs:3.0,0);
  \node[font=\sffamily\footnotesize, anchor=south, fill=white, inner sep=1pt] at (axis cs:1.3,3.05) {① $t=0.040$};
  \node[font=\sffamily\footnotesize, anchor=east, fill=white, inner sep=1pt] at (axis cs:1.97,2.3) {② $t=0.080$};
  \node[font=\sffamily\footnotesize, anchor=north, fill=white, inner sep=1pt] at (axis cs:1.1,-1.08) {③ $t=0.120$：$-1$};
  \node[font=\sffamily\footnotesize, anchor=west, inner sep=1pt] at (axis cs:2.53,2.2) {③：2};
  \node[font=\sffamily\footnotesize, anchor=north, inner sep=1pt] at (axis cs:2.55,3.55) {弦2（重い）};
\end{axis}
\end{tikzpicture}
\bkfigcaption{表計算で動かした弦の形。①（細い実線）高さ\qty{3}{cm}の山が右へ進む。②（破線）つなぎ目（$x=\qty{2.0}{m}$）を通るところ。③（太い実線）はね返った山は反対側に高さ\qty{1}{cm}、通りぬけた山は高さ\qty{2}{cm}で、はばが半分。}\label{ph:波-7:fig-sheet}
\end{bkfigure}

\noindent 結果は、はね返った山が反対側に\qty{1}{cm}、通りぬけた山が\qty{2}{cm}でした。④の式で $r=-\tfrac13$、$\tau=\tfrac23$ なので、$3\times(-\tfrac13)=-1$、$3\times\tfrac23=2$ と、ぴったり合います。通りぬけた山のはばが半分なのは、右の弦の速さが半分だからです（同じ時間に、半分の距離しか進まない）。右の弦を軽く（$Z_2=\tfrac12Z_1$）すると、はね返った山は同じ側に\qty{1}{cm}、通りぬけた山は\qty{4}{cm}になります（計算してみよう）。

\begin{seikaku}
つなぎ目の式は、つなぎ目の点に、左の弦の半マス分と右の弦の半マス分の質量をもたせ、両側の弦の張力で引かせて、波-5 ⑧と同じ1歩の式を作ると出てきます（章末問題）。この式は、$Z_1$ と $Z_2$ だけで書けているのがよい所で、④の決まり1・決まり2を、表計算の中で守らせていることになります。
\end{seikaku}

\phboxhead{はてな？の答え}
\label{ph:波-7:hatena}\bkhatenaref の答えです。プローブを乾いたおなかに当てると、どんなにぴったり押しつけたつもりでも、皮膚の細かい凹凸のあいだに、薄い空気の層が残ります。空気（$Z=413$）から体（$Z=1.63\times10^6$）へ入るときは、比が約4000倍もあるので、
\[r=\frac{413-1.63\times10^6}{413+1.63\times10^6}\approx-0.9995,\qquad r^2\approx0.999\]
で、\textbf{エネルギーの99.9\,\%がはね返り}、体の中へ入るのは0.1\,\% だけです（実際は、プローブから空気へ出るときにも、ほとんどがはね返ります）。しかも、体の中からはね返ってきた音が、もう一度この境目を通って戻るので、さらに0.1\,\% になります。これでは、絵が作れません。

ゼリーは、ほとんど水でできていて、$Z$ が約 $1.5\times10^6$ と、体に近い値です。すきまの空気を追い出してゼリーで満たすと、ゼリーから体へは
\[r=\frac{1.5\times10^6-1.63\times10^6}{1.5\times10^6+1.63\times10^6}\approx-0.042,\qquad r^2\approx0.0017\]
で、はね返るのは0.17\,\% だけ、99.8\,\% が体の中へ入ります。\textbf{ゼリーは、インピーダンスの大きく違う空気を、すきまから追い出すためにぬる}のです。

\begin{phmodel}[2本の弦・2つの物の境目]
\phignore{境目の結び目の質量と大きさ。媒質の中での吸収。}
\phassume{波は境目に垂直に入る（音なら、平らな面に平らな波が垂直に入る）。ゆれは小さく、$Z$ はゆれの大きさによらない。}
\phlimit{斜めに入るとき（屈折が起きる：波-8 ①）。境目が曲がっていたり、でこぼこしていたりするとき。境目の層の厚さが、波長と同じくらいのとき（⑧）。}
\end{phmodel}

\phkokomade

\Needspace{16\baselineskip}\phstage{shiki}{波の式と境界条件から、反射を導く}{5}

\bksubsection{5}{境界条件から反射係数を導く}
\label{ph:波-7:sec5}
④では、「力＝$Z$×速さ」（③）を使って $r$ と $\tau$ を求めました。ここでは、波動方程式の答え（波-5 ④）から出発して、同じ式を導き直します。そうすると、弦のはしに、ばねやおもりやダンパー（ゆれを止める筒）をつけたときにも、同じ考えで反射が計算できるようになります。

\phboxhead{2つの弦の波を式で書く}
つなぎ目を $x=0$ にとり、左（$x<0$）を弦1、右（$x>0$）を弦2とします。張力はどちらも $F_{\mathrm T}$（つなぎ目の左右の引き合いがつり合うので同じ）、線密度は $\mu_1$・$\mu_2$、波の速さは $c_1=\sqrt{F_{\mathrm T}/\mu_1}$・$c_2=\sqrt{F_{\mathrm T}/\mu_2}$ です。どちらの弦でも、波動方程式
\[\frac{\partial^2y}{\partial t^2}=c^2\,\frac{\partial^2y}{\partial x^2}\]
が成り立ちます（$c$ は、それぞれの弦の $c_1$・$c_2$）。来た波・はね返った波・通りぬけた波を、複素数で書きます（波-5 ⑨。本物の波は実部）。
\[y_1=Ae^{j(\omega t-k_1x)}+rAe^{j(\omega t+k_1x)}\quad(x<0),\qquad y_2=\tau Ae^{j(\omega t-k_2x)}\quad(x>0)\]
$e^{j(\omega t+k_1x)}$ は、左へ進む波です（$x$ が減る向きに形が動く）。3つの波は、同じ $\omega$ をもちます。つなぎ目は1つの点なので、両側を同じリズムでゆらすからです。波数は、それぞれの弦で $k_1=\omega/c_1$、$k_2=\omega/c_2$ と違います（波長が違う）。

\phboxhead{境目で守る2つの条件（境界条件）}
\begin{enumerate}[label={\bkcircled{\arabic*}}, leftmargin=2.4\zw, itemsep=2pt, topsep=2pt]
\item \textbf{弦がつながっている}：$x=0$ で、$y_1=y_2$。
\item \textbf{結び目の力がつり合う}：結び目の質量は0なので、左右の張力のたての成分がつり合う。たての成分は $F_{\mathrm T}\times$（傾き）なので、$x=0$ で $F_{\mathrm T}\,\dfrac{\partial y_1}{\partial x}=F_{\mathrm T}\,\dfrac{\partial y_2}{\partial x}$。
\end{enumerate}
このような、境目で守る条件を\textbf{境界条件}といいます（8-13③の「両はしの決まり」と同じ考え）。\bkidx{きょうかいじょうけん}{境界条件}

\phboxhead{条件を式にする}
$x=0$ を入れると、$e^{\pm jk_1x}$ も $e^{-jk_2x}$ も1なので、①は $Ae^{j\omega t}$ でわって
\[1+r=\tau\]
②では、$x$ で微分すると $e^{-jkx}$ の前に $-jk$、$e^{jkx}$ の前に $jk$ が出ます（6-10⑩）。$x=0$ で、
\[F_{\mathrm T}(-jk_1+jk_1r)Ae^{j\omega t}=F_{\mathrm T}(-jk_2\tau)Ae^{j\omega t}\]
両辺を $-jAe^{j\omega t}\omega$ でわり、$F_{\mathrm T}k/\omega=F_{\mathrm T}/c=Z$（③）を使うと、
\[Z_1(1-r)=Z_2\,\tau\]
です。これは、④の決まり2と同じ式です。2つを解くと、
\pheq{r=\frac{Z_1-Z_2}{Z_1+Z_2}\qquad\tau=\frac{2Z_1}{Z_1+Z_2}}{単位なし}
が導けました。④の答えと同じです。エネルギーについても、
\[r^2+\frac{Z_2}{Z_1}\tau^2=\frac{(Z_1-Z_2)^2+4Z_1Z_2}{(Z_1+Z_2)^2}=\frac{(Z_1+Z_2)^2}{(Z_1+Z_2)^2}=1\]
で、はね返る分と通りぬける分の和は、来た分とちょうど等しくなっています。

\phboxhead{軽いひもから重いひもへ：山と谷が入れかわる}
\label{ph:波-7:isou}軽いひもから重いひもへ進むとき（$Z_1<Z_2$）は、$r$ が負になります。$r=-|r|$ と書くと、$-1=e^{j\pi}$（6-10⑥）なので、はね返った波は
\[rAe^{j(\omega t+k_1x)}=|r|\,Ae^{j(\omega t+k_1x+\pi)}\]
です。山が谷に、谷が山に入れかわってはね返ります。これを「はね返るとき\textbf{位相が $\pi$（180°）ずれる}」といいます。\bkidx{いそうのずれ}{位相のずれ（反射）}重いひもから軽いひもへ（$Z_1>Z_2$）なら $r>0$ で、位相はずれません。通りぬける波は、$\tau=2Z_1/(Z_1+Z_2)$ がいつも正なので、どちらの向きでも位相はずれません。光の薄い膜の干渉（波-13 ④）では、この「どちらの面で位相が $\pi$ ずれるか」が答えを決めます（⑨）。

\Needspace{14\baselineskip}
\begin{phflow}[反射係数と固定端・自由端]
\phflowrow{境目で $y$ と $F_{\mathrm T}\dfrac{\partial y}{\partial x}$ がつながる}{平らな波が\\垂直に入る}{$r=\dfrac{Z_1-Z_2}{Z_1+Z_2}$}
\phflowrow{$r=\dfrac{Z_1-Z_2}{Z_1+Z_2}$}{$Z_2\to\infty$\\（動けない端）}{$r=-1$（固定端）}
\phflowrow{$r=\dfrac{Z_1-Z_2}{Z_1+Z_2}$}{$Z_2\to0$\\（自由に動ける端）}{$r=+1$（自由端）}
\end{phflow}
\label{ph:波-7:koteitan}壁にしっかり結んだはしは、どんなに強く押しても動かないので、$Z_2$ が限りなく大きい物とみなせます。分子と分母を $Z_2$ でわると $r=(Z_1/Z_2-1)/(Z_1/Z_2+1)\to-1$ で、山は同じ高さの谷になってはね返ります。なめらかな棒にはめた軽い輪のように、少しの力でいくらでも動くはしは $Z_2\to0$ で、$r\to+1$、山は山のままはね返ります。波-6 ①の固定端と自由端は、この式の両はしだったのです。

\phboxhead{はしに物をつける}
弦1のはしに、何かの装置をつけます。装置が、はしの速さ $u$ に対して $F=Z_{\text{端}}u$ の力で押し返すなら、装置は「インピーダンス $Z_{\text{端}}$ の弦2」と同じ働きをします。だから、$r=(Z_1-Z_{\text{端}})/(Z_1+Z_{\text{端}})$ です。
\begin{itemize}[leftmargin=1.2\zw, itemsep=2pt, topsep=2pt, label={\color{bkTheme}●}]
\item \textbf{ダンパー}（速さに比例して止める力 $F=bu$、$b$ は係数）：$Z_{\text{端}}=b$。$b=Z_1$ にすると $r=0$ で、まったくはね返りません（\textbf{はね返らない端}、無反射終端。例題11）。来た波のエネルギーは、全部ダンパーの中で熱になります。
\item \textbf{ばね}（ばね定数 $k_{\mathrm s}$、力 $F=k_{\mathrm s}y$）：$y=Ye^{j\omega t}$ とすると、速さは $u=j\omega Y e^{j\omega t}$（6-10⑩）なので、
\[F=k_{\mathrm s}y=\frac{k_{\mathrm s}}{j\omega}\,u,\qquad Z_{\text{端}}=\frac{k_{\mathrm s}}{j\omega}=-j\,\frac{k_{\mathrm s}}{\omega}\]
と、インピーダンスが\textbf{複素数}になります。
\end{itemize}

ばねのはしの反射係数は、分子と分母に $j$ がある複素数です。
\[r=\frac{Z_1+j\,k_{\mathrm s}/\omega}{Z_1-j\,k_{\mathrm s}/\omega}\]
分子と分母は、複素平面で、実軸について折り返した位置にある数（大きさが同じで、角が反対）です。わり算は「大きさはわり算、角はひき算」（6-9⑦）なので、$|r|=1$、$r$ の角は分子の角の2倍の $2\arctan\bigl(k_{\mathrm s}/(\omega Z_1)\bigr)$ です。\textbf{ばねは、エネルギーを全部はね返し（$|r|=1$）、はね返る波の山の時刻だけをずらす}のです。ばねはエネルギーをためて返すだけで、熱に変えないからです。

\begin{bkcenter}
\begin{bktable}{colspec={Q[c,m,wd=26mm] Q[c,m,wd=20mm] Q[c,m,wd=24mm] Q[l,m,wd=42mm]}, rowsep=1.5pt, row{1}={bg=bkTheme!10, font=\sffamily\small}}
振動数 $f$〔Hz〕 & $|r|$ & $r$ の角 & はね返り方 \\
0.1 & 1 & 145° & 固定端（180°）に近い \\
1 & 1 & 35° & あいだ \\
10 & 1 & 3.6° & 自由端（0°）に近い \\
\end{bktable}
\bktabcaption{弦（$Z_1=\qty{4.0}{N.s/m}$）のはしに、ばね定数\qty{8}{N/m}のばねをつけたときの反射係数。ゆっくりゆすると、ばねはかたい壁のように、速くゆすると、やわらかい輪のようにふるまう。}\label{ph:波-7:tab-bane}
\end{bkcenter}

\begin{seikaku}
音で圧力の反射係数 $r_{\text{押}}=(Z_2-Z_1)/(Z_1+Z_2)$（④の2つの数え方）を使うと、空気から水へ進む音では $r_{\text{押}}$ はほぼ $+1$ で、水面では圧力が2倍近くになり、水面の速さはほぼ0です。波-6 ④の「閉じた管のはしは、圧力の腹」と同じことです。どちらの符号で書いても、はね返るエネルギーの割合 $r^2$ は同じです。
\end{seikaku}

\bksubsection{6}{分散のある波：エネルギーは群速度で進む}
\label{ph:波-7:sec6}
①では、エネルギーは波といっしょに速さ $c$ で運ばれる、としました。これは、弦の波や空気中の音のように、どの振動数の波も同じ速さで進むときの話です。深い海の波は、波長が長いほど速く進みます。このように、波の速さが振動数（波長）によって違うことを\textbf{分散}といい、$\omega$ と $k$ の関係を\textbf{分散関係}といいます。波-5 ⑦で、振動数の少し違う2つの波の和（4-9⑦のうなり）を考え、うなりのつつみ（ふくらみ）の形が\textbf{群速度} $c_{\mathrm g}=\dd\omega/\dd k$ で進むことを学びました。\bkidx{ぐんそくど}{群速度}

波が運ぶエネルギーは、山の1つ1つではなく、つつみの形といっしょに進みます。エネルギーがあるのは、ゆれの大きい所（つつみのふくらんだ所）だからです。分散のある波では、エネルギーの流れは
\[P=\overline E\,c_{\mathrm g}\]
となり、①の $c$ を $c_{\mathrm g}$ に置きかえます。分散のない波（$\omega=ck$）では $c_{\mathrm g}=c$ で、①と同じです。

\phboxhead{深い海の波}
深い海（深さが波長の半分より深い所）の波は、流体力学で導かれる分散関係 $\omega=\sqrt{gk}$ に従います（この式の導き方は専門の本で）。山の進む速さ（\textbf{位相速度}）$c=\omega/k$ と群速度は、
\[c=\frac{\omega}{k}=\sqrt{\frac{g}{k}},\qquad c_{\mathrm g}=\frac{\dd\omega}{\dd k}=\frac{\dd}{\dd k}\bigl(\sqrt g\,k^{1/2}\bigr)=\frac12\sqrt{\frac{g}{k}}=\frac{c}{2}\]
で、エネルギーは山の半分の速さでしか進みません（例題8）。

\begin{bkfigure}
\begin{tikzpicture}
\begin{axis}[bkplot, name=g1, width=124mm, height=30mm, xmin=0, xmax=40, ymin=-1.3, ymax=1.3, xtick=\empty, ytick=\empty, xlabel={$x$}, axis lines=middle]
  \addplot[bkcurve, domain=0:40, samples=500] {exp(-((x-13)/5)^2)*cos(deg(2*pi*x/2.5))};
  \addplot[bkOrange, line width=1pt, dash pattern=on 3pt off 1.5pt, domain=0:40, samples=120] {exp(-((x-13)/5)^2)};
  \fill[bkRed] (axis cs:12.5,0.99) circle (1.6pt);
  \node[font=\sffamily\footnotesize, anchor=west, fill=white, inner sep=1pt] at (axis cs:24,0.9) {はじめ（●は1つの山）};
\end{axis}
\begin{axis}[bkplot, at={(g1.south)}, anchor=north, yshift=-6mm, width=124mm, height=30mm, xmin=0, xmax=40, ymin=-1.3, ymax=1.3, xtick=\empty, ytick=\empty, xlabel={$x$}, axis lines=middle]
  \addplot[bkcurve, domain=0:40, samples=500] {exp(-((x-21)/5)^2)*cos(deg(2*pi*(x-16)/2.5))};
  \addplot[bkOrange, line width=1pt, dash pattern=on 3pt off 1.5pt, domain=0:40, samples=120] {exp(-((x-21)/5)^2)};
  \fill[bkRed] (axis cs:28.5,0.105) circle (1.6pt);
  \draw[bkGray, line width=0.4pt, dashed] (axis cs:13,-1.3) -- (axis cs:13,1.3);
  \node[font=\sffamily\footnotesize, anchor=west, fill=white, inner sep=1pt] at (axis cs:31,1.1) {しばらくあと};
  \path (axis cs:0,-1.9);
\end{axis}
\end{tikzpicture}
\bkfigcaption{深い海の波の群れ（実線）と、そのつつみ（破線）。つつみ（エネルギー）は位置13から21へ進んだが、●の山は12.5から28.5へと、2倍の距離を進んで、群れの前のはしで消えかけている。下の破線は、はじめのつつみのまん中の位置。群れの後ろでは、新しい山が生まれる。}\label{ph:波-7:fig-gun}
\end{bkfigure}

\begin{phmijika}[遠くの台風が送るうねり]
夏、晴れているのに、海岸に大きなゆったりした波（うねり）が来ることがあります。遠くの台風の所で起きた波が、何日もかけて届いたものです。群速度 $c_{\mathrm g}=c/2$ は、波長が長い（周期が長い）ほど大きいので、\textbf{周期の長いうねりが先に}着き、だんだん周期の短い波が着くようになります。うねりの周期の変わり方から、台風までの距離の見当がつきます（章末問題）。
\end{phmijika}

\bksubsection{7}{弱まる波：広がりと吸収}
\label{ph:波-7:sec7}
波が進むにつれて弱くなる理由は、2つあります。1つは、②で見た\textbf{広がり}、もう1つは、媒質がエネルギーを熱に変えてしまう\textbf{吸収}です。

\phboxhead{球面波の式}
②では、エネルギーの数え方から $I\propto1/R^2$ を出しました。ここでは、波の式から確かめます。音の圧力 $p$ は、波-5 ④の1次元の波動方程式を3次元にした
\[\frac{\partial^2p}{\partial t^2}=c^2\nabla^2p\]
をみたします（$\nabla^2$ はラプラシアン：8-12）。小さな音源から四方へ同じように広がる波では、$p$ は音源からの距離 $R$ と時刻 $t$ だけで決まるので、8-12⑦の球のラプラシアンが使えます。
\[\nabla^2p=\frac{1}{R^2}\,\frac{\partial}{\partial R}\Bigl(R^2\,\frac{\partial p}{\partial R}\Bigr)=\frac{\partial^2p}{\partial R^2}+\frac{2}{R}\,\frac{\partial p}{\partial R}\]
ここで、$Rp$ を $R$ で2回微分すると（積の微分：5-6）、
\[\frac{\partial^2(Rp)}{\partial R^2}=R\,\frac{\partial^2p}{\partial R^2}+2\,\frac{\partial p}{\partial R}\qquad\text{なので}\qquad\nabla^2p=\frac1R\,\frac{\partial^2(Rp)}{\partial R^2}\]
と書けます。波動方程式の両辺に $R$ をかけると、
\[\frac{\partial^2(Rp)}{\partial t^2}=c^2\,\frac{\partial^2(Rp)}{\partial R^2}\]
で、$Rp$ が1次元の波動方程式をみたします。外へ進む波の答えは、ダランベールの解（8-13⑨）と同じ形の $Rp=F(R-ct)$ です。
\pheq{p=\frac{F(R-ct)}{R}}{Pa ＝ Pa·m÷m}
\label{ph:波-7:kyumenha}形 $F$ は、そのまま外へ進みますが、高さは $1/R$ に比例して低くなります。強さは $p^2$ に比例するので $1/R^2$ に比例し、②の逆2乗の法則と合います。この式を\textbf{球面波}の式といいます。\bkidx{きゅうめんは}{球面波}

\phboxhead{吸収：同じ距離ごとに同じ割合で減る}
空気や水や体の中では、音の圧力の変化で、物の中にわずかな粘りによるこすれや、熱の出入りが起き、波のエネルギーの一部が熱に変わります。1 m 進むごとに、振幅が同じ割合で減るので、平面の波の振幅は
\pheq{A(x)=A_0\,e^{-\alpha x}}{1/m の $\alpha$ と m の $x$ で、肩は単位なし}
のように指数関数で減ります（5-14）。$\alpha$〔1/m〕を\textbf{吸収の定数}（減衰定数）といいます。\bkidx{きゅうしゅう}{吸収（波）}強さは振幅の2乗なので $e^{-2\alpha x}$ で減ります。工学では、弱まり方を dB で書きます。振幅が $e^{-\alpha x}$ 倍になると、
\[20\log_{10}e^{\alpha x}=20\,(\log_{10}e)\,\alpha x\approx8.69\,\alpha x\ \text{〔dB〕}\]
なので、「1 m で何 dB」と書けば、距離に比例して dB をたしていけばよいのです（3-12⑦の「かけ算がたし算になる」）。

\phboxhead{高い音ほど、よく吸収される}
吸収の定数は、ふつう振動数が高いほど大きくなります。遠くの雷が「ゴロゴロ」と低く聞こえるのは、何 km も進むうちに、高い音が空気に吸収されて、低い音だけが残るからです。体の軟らかい組織では、吸収はおよそ振動数に比例し、1 cm・1 MHz あたり約\qty{0.5}{dB}です。
\[\text{弱まり〔dB〕}\approx0.5\times f\text{〔MHz〕}\times\text{進む距離〔cm〕}\]
超音波の検査では、音は体の中を行って戻るので、進む距離は深さの2倍です（例題9）。

\begin{phbaisuru}{超音波の振動数 $f$}
吸収の dB も2倍になります。いっぽう、波長は半分になり、細かい物まで見分けられます（波長より小さい物は見えにくい）。深い所は低い振動数（\qty{3.5}{MHz}くらい）、皮膚に近い所は高い振動数（\qty{10}{MHz}くらい）で調べるのは、このためです。
\par\smallskip\noindent\begin{tikzpicture}
\begin{axis}[bkplot, width=96mm, height=44mm, xmin=0, xmax=17, ymin=0, ymax=165, xtick={5,10,15}, ytick={50,100,150}, xlabel={深さ（cm）}, ylabel={弱まり（dB）}]
  \addplot[bkcurve, domain=0:16] {0.5*3.5*2*x};
  \addplot[bkOrange, line width=1.6pt, dash pattern=on 4pt off 2pt, domain=0:16] {0.5*5*2*x};
  \addplot[bkRed!80!black, line width=1.6pt, dash pattern=on 1.2pt off 1.4pt, domain=0:15] {0.5*10*2*x};
  \node[font=\sffamily\footnotesize, anchor=west, fill=white, inner sep=1pt] at (axis cs:16.2,56) {3.5 MHz};
  \node[font=\sffamily\footnotesize, anchor=west, fill=white, inner sep=1pt] at (axis cs:16.2,80) {5 MHz};
  \node[font=\sffamily\footnotesize, anchor=south east, fill=white, inner sep=1pt] at (axis cs:14.5,150) {10 MHz};
\end{axis}
\end{tikzpicture}
\par\noindent（行って戻るときの弱まり。実線 3.5 MHz、破線 5 MHz、点線 10 MHz。）
\end{phbaisuru}

\begin{phmada}[体の中の超音波の吸収は、なぜ振動数にほぼ比例する？]
水だけなら、音の吸収は振動数の2乗に比例することが、粘りの理論から説明できます。ところが、体の軟らかい組織や多くのやわらかい材料では、吸収が振動数の1乗に近い形（多くの組織で1乗から1.7乗くらい）で増えることが、たくさんの測定でわかっています。たんぱく質の大きな分子や細胞のつくりなど、大きさの違うたくさんの原因が重なって、この形になると考えられていますが、組織のつくりから吸収の大きさと振動数への依存を計算して説明し切ることは、まだできていません。この性質をうまく表す新しい式（分数回の微分を使う式など）も研究されています。
\end{phmada}

\bksubsection{8}{はね返りを減らす：整合}
\label{ph:波-7:sec8}
インピーダンスの違う2つの物のあいだで、はね返りを減らしてエネルギーをよく通すくふうを、\textbf{整合}（マッチング）といいます。\bkidx{せいごう}{整合}ゼリーは、空気を追い出す整合の1つでした。

\phboxhead{4分の1波長の層}
超音波のプローブの中では、圧電セラミックス（$Z_1\approx3\times10^7$）が電気で音を作ります。ゼリーや体（$Z_2\approx1.5\times10^6$）へ直接出すと、$Z$ が20倍違うので、エネルギーの約8割がはね返ってしまいます。そこで、あいだに、インピーダンス $Z_{\mathrm m}$、厚さが層の中の波長 $\lambda_{\mathrm m}$ の $\tfrac14$ の層をはります（\textbf{整合層}）。
\begin{bkfigure}
\begin{tikzpicture}[x=1mm, y=1mm]
  \fill[bkNavy!25] (0,0) rectangle (30,30);
  \fill[bkYellow!35] (30,0) rectangle (46,30);
  \fill[bkGreen!25] (46,0) rectangle (80,30);
  \draw[bkInk, line width=0.5pt] (30,0) -- (30,30) (46,0) -- (46,30);
  \node[font=\sffamily\footnotesize, align=center] at (15,4) {$Z_1$};
  \node[font=\sffamily\footnotesize, align=center] at (38,4) {$Z_{\mathrm m}$};
  \node[font=\sffamily\footnotesize, align=center] at (63,4) {$Z_2$};
  \draw[bkBlue!80!black, line width=1pt, -{Stealth[length=2.2mm]}] (4,24) -- (29,24);
  \node[font=\sffamily\footnotesize, anchor=south] at (16,24.6) {来た波};
  \draw[bkRed!80!black, line width=1pt, -{Stealth[length=2.2mm]}] (29,21) -- (6,21);
  \node[font=\sffamily\footnotesize, anchor=north] at (17,20.3) {ア（はじめの面）};
  \draw[bkOrange!80!black, line width=1pt] (31,24) -- (45,24);
  \draw[bkOrange!80!black, line width=1pt, -{Stealth[length=2.2mm]}] (45,24) -- (45,14) -- (6,14);
  \node[font=\sffamily\footnotesize, anchor=north] at (17,13.3) {イ（2つめの面）};
  \draw[bkInk, line width=0.5pt, {Stealth[length=2pt]}-{Stealth[length=2pt]}] (30,32) -- (46,32);
  \node[font=\sffamily\footnotesize, anchor=south] at (38,32.5) {$\lambda_{\mathrm m}/4$};
  \node[anchor=west, font=\small, text width=40mm, align=left] at (84,16) {%
  イは、層の中を行って戻る分 $\lambda_{\mathrm m}/2$ だけ遠回り。\par\vspace{2pt}
  半波長ずれると、山と谷が重なって、アとイは打ち消し合う。};
\end{tikzpicture}
\bkfigcaption{4分の1波長の整合層。はじめの面ではね返る波アと、2つめの面ではね返る波イが打ち消し合う。}\label{ph:波-7:fig-seigou}
\end{bkfigure}

\noindent はね返る波は2つあります。層のはじめの面ではね返る波アと、層を通りぬけて2つめの面ではね返り、戻ってくる波イです。イは、層の中を行って戻る分、$2\times\tfrac14\lambda_{\mathrm m}=\tfrac12\lambda_{\mathrm m}$ だけ遠回りするので、アと半波長ずれて、山と谷が重なります。アとイの高さが同じなら、ぴったり打ち消し合って、はね返りがなくなります。アの反射係数は $(Z_1-Z_{\mathrm m})/(Z_1+Z_{\mathrm m})$、イは $(Z_{\mathrm m}-Z_2)/(Z_{\mathrm m}+Z_2)$ です。この2つが等しいのは、$Z_1/Z_{\mathrm m}=Z_{\mathrm m}/Z_2$ のときなので、
\pheq{Z_{\mathrm m}=\sqrt{Z_1Z_2}\qquad\text{厚さ}=\frac{\lambda_{\mathrm m}}{4}}{Pa·s/m ＝ $\sqrt{(\text{Pa·s/m})^2}$}
\label{ph:波-7:seigou}です（2つの $Z$ の\textbf{相乗平均}）。ここでは、イが層の境目でさらに何回もはね返る分を無視しましたが、それも入れて正しく計算しても、設計した振動数では、はね返りはちょうど0になります（専門の本で）。振動数がずれると、半波長の打ち消しがずれて、少しはね返ります。

\begin{phmijika}[耳の中の整合：中耳]
耳では、空気の音が、鼓膜から、小さな3つの骨（耳小骨）を通って、液体の入った内耳へ伝わります。空気から液体（水に近い）へ直接入ると、エネルギーは約0.1\,\%（$-30$\,dB）しか入りません（例題5）。中耳では、広い鼓膜で受けた力を、面積がその17〜20分の1ほどの小さなあぶみ骨の底に集め、骨のてこ（約1.3倍）も使って、計算の上では圧力を20倍あまり（約\qty{27}{dB}）にしてから内耳へ伝えます（面積とてこの比から出す目安。実際の働きは音の高さで変わり、もう少し小さい）。小さな面積を大きな力で押す、つまりインピーダンスの大きな液体に合わせる整合のしかけです。
\end{phmijika}

\phboxhead{ラッパとメガホン}
ラッパやメガホンは、口元の細い管から、大きく開いた口まで、少しずつ広がっています。管の中の音では、管全体を1つの通り道と見て「圧力÷（1秒に動く空気の体積）」を考えると、これは $\rho c/S$（$S$ は断面積）で、太い所ほど小さくなります。管をゆっくり広げると、この値が少しずつ変わるので、はね返りを少なくして、音を外の広い空気へ出せます。インピーダンスの変わり方が、波長に比べてゆるやかなほど、はね返りは少なくなります。

\phboxhead{電気の線の整合}
テレビのアンテナ線（同軸ケーブル）を伝わる電気の波でも、つないだ先のインピーダンスが違うと、はね返りが起きます。電圧（押すもの）の反射係数は、④の $r_{\text{押}}$ と同じく $(Z_{\text{先}}-Z_{\text{線}})/(Z_{\text{先}}+Z_{\text{線}})$ です。はね返った波は、画面の「ゴースト」（二重の像）や、通信の誤りのもとになります。だから、ケーブルとつなぐ機器は、同じインピーダンス（テレビは\qty{75}{\ohm}、無線機は\qty{50}{\ohm}が多い）にそろえます（例題7。詳しくは電-13）。

\bksubsection{9}{光の反射率}
\label{ph:波-7:sec9}
光も波です（電磁波。電-15）。光がガラスや水の表面に垂直に当たるときも、インピーダンスの違いで一部がはね返ります。磁石につかない透明な物では、光にとってのインピーダンスは、屈折率 $n$（光の進みにくさ。波-12 ③）に反比例します。④の式で $Z\propto1/n$ とすると、はね返る割合は
\pheq{r^2=\Bigl(\frac{n_1-n_2}{n_1+n_2}\Bigr)^2}{単位なし}
\label{ph:波-7:hikari}になります。空気（$n_1=1.00$）からガラス（$n_2=1.5$）へ垂直に入ると、$(0.5/2.5)^2=0.04$ で、\textbf{4\,\%} がはね返ります（例題12）。水（$n=1.33$）では約2\,\%、ダイヤモンド（$n=2.42$）では約17\,\% で、ダイヤモンドがきらきら光る理由の1つです。

\begin{phmijika}[夜の窓ガラス]
昼間、窓の外はよく見えます。夜、明るい部屋から暗い外を見ると、窓に自分の顔や部屋の中がはっきり写ります。窓ガラスは、表と裏の2つの面で、合わせて約8\,\% の光をはね返しています。昼は、外から入ってくる光がずっと強いので、はね返った8\,\% は目立ちません。夜は外が暗いので、部屋の光のはね返りだけが見えるのです。ガラスが変わったのではなく、比べる光の強さが変わったのです。
\end{phmijika}

\noindent 光では、屈折率の大きい物の表面ではね返るとき、電場（光の「押すもの」の側）の山が谷になります（④の $r_{\text{押}}=(Z_2-Z_1)/(Z_1+Z_2)$ で、$Z\propto1/n$ なので $n_2>n_1$ なら $Z_2<Z_1$ となり、負になる）。このことは、シャボン玉の色（薄い膜の干渉）を考えるときに使います（波-13 ④）。メガネやカメラのレンズの\textbf{反射防止膜}は、⑧の4分の1波長の層を光で作ったもので、屈折率が $\sqrt{1.5}\approx1.22$ に近い物（フッ化マグネシウム、約1.38）の薄い膜をつけて、はね返りを4\,\% から1\,\% 台に減らしています（波-13）。

\begin{phsaki}{波の反射と伝送}
この章では、平らな波が境目に垂直に入る場合だけを考えました。斜めに入ると、屈折が起き（波-8 ①）、はね返る割合も角度で変わります（光では、ある角度ではね返りが0になる偏光：波-13）。電気の線の反射は「分布定数回路」として、アンテナや高速のデジタル回路の設計に欠かせません（電-13）。光が電磁波であることから $Z\propto1/n$ を導くのは電-15 です。また、ふつうの物にはない $Z$ を、小さな部品をたくさん並べて作る「音響メタマテリアル」の研究も進んでいます。
\end{phsaki}

\begin{senmonsho}[音響工学の専門書の書き方]
\begin{bkhonbun}
媒質 1 から媒質 2 へ平面波が垂直入射するとき，音圧反射係数 $R_p$，音圧透過係数 $T_p$，強度透過率 $\tau_I$ は
\[R_p=\frac{Z_2-Z_1}{Z_2+Z_1},\qquad T_p=\frac{2Z_2}{Z_2+Z_1},\qquad \tau_I=\frac{4Z_1Z_2}{(Z_1+Z_2)^2}\]
で与えられる．ここで $Z=\rho c$ は特性音響インピーダンス〔Rayl〕である．
\end{bkhonbun}
\bkyaku{音響の本は、ふつう圧力（押すもの）で反射を書くので、この章の $r$ と符号が逆の $R_p=-r$。圧力の透過係数 $T_p$ は $(Z_2/Z_1)\tau$。Rayl（レイル）は Pa·s/m のこと（イギリスの物理学者レイリーの名前から）。$1.5\times10^6$ Pa·s/m を 1.5 MRayl と書く。強度透過率の $\tau_I$ は、この章の「通りぬける割合」。}
\end{senmonsho}

\begin{taisetsu}
\def\tsl#1{\makebox[7.6\zw][l]{\sffamily\bfseries\small\color{\bkTT}#1}}
\tsl{波のエネルギー}$\overline E=\dfrac12\mu\omega^2A^2$〔J/m〕、$P=\overline E\,c=\dfrac12Z\omega^2A^2$\\[3pt]
\tsl{強さ}$I=\dfrac{P}{4\pi R^2}$（点の音源）、音では $I=\dfrac{p_{\mathrm e}^{\,2}}{\rho c}$\\[3pt]
\tsl{インピーダンス}$Z=\dfrac{\text{押すもの}}{\text{流れるもの}}$：弦 $\sqrt{F_{\mathrm T}\mu}=\mu c$、音 $\rho c$\\[3pt]
\tsl{反射と透過}$r=\dfrac{Z_1-Z_2}{Z_1+Z_2}$、$\tau=\dfrac{2Z_1}{Z_1+Z_2}$、割合 $r^2$ と $\dfrac{4Z_1Z_2}{(Z_1+Z_2)^2}$\\[3pt]
\tsl{弱まり}球面波 $p=\dfrac{F(R-ct)}{R}$、吸収 $e^{-\alpha x}$（$8.69\alpha$ dB/m）\\[3pt]
\tsl{整合}$Z_{\mathrm m}=\sqrt{Z_1Z_2}$、厚さ $\dfrac{\lambda_{\mathrm m}}{4}$。光は $\Bigl(\dfrac{n_1-n_2}{n_1+n_2}\Bigr)^2$
\end{taisetsu}

\begin{kotoba}
\small
\bkword{強さ}{波の進む向きに直角な面 1 m$^2$ を、1秒に通るエネルギー。単位 W/m$^2$。}
\bkword{インピーダンス}{媒質を動かすのに要る「押すもの」÷「流れるもの」。弦は力÷速さ、音は圧力÷速さ。}
\bkword{反射係数・透過係数}{はね返る波・通りぬける波の高さ（速さ）の、来た波に対する比。}
\bkword{境界条件}{境目で守る条件。弦では、つながっていることと、力がつり合うこと。}
\bkword{群速度}{波の群れ（つつみ）とエネルギーが進む速さ $\dd\omega/\dd k$。}
\bkword{整合}{インピーダンスをそろえて、はね返りを減らすこと。}
\end{kotoba}

\bkkigocards{\bkkigocard{$Z$}{ゼット}{インピーダンス。弦は N·s/m、音は Pa·s/m、電気は Ω。第1巻の断面係数 $Z$ とは別の量（この巻では断面係数を使わない）}{impedance}{波-7}}
{\bkkigocard{$I$}{アイ}{波の強さ〔W/m$^2$〕。第1部（波-1・3）の慣性モーメント $I$、波-6 ⑨の断面二次モーメント $I$ とは別の量}{intensity}{波-7}}
\bkkigocards{\bkkigocard{$r$}{アール}{振幅（変位・速さ）の反射係数。音響の本の音圧反射係数は $-r$。はね返るエネルギーの割合は $r^2$}{reflection coefficient}{波-7}}{}

%% ======================== 例題 ========================
\Needspace{10\baselineskip}\bksection{reidai}{例題でたしかめよう}
\label{ph:波-7:ex1}
\begin{phreidai}[ゼリーをぬる・ぬらない（はてな？）]
超音波の検査で、プローブと体のあいだに、（1）空気の薄い層があるとき、（2）ゼリーで満たしたとき、境目（空気→体、ゼリー→体）ではね返るエネルギーの割合と、体へ入る割合を求めましょう。$Z$ は、空気413、ゼリー $1.5\times10^6$、体 $1.63\times10^6$ Pa·s/m とします。
\phbamen
音が、インピーダンスの違う2つの物の境目に、垂直に入る場面です。
\phmushi
プローブの側の境目（プローブ→空気・ゼリー）は考えません。境目は平らで、吸収はないとします。
\phzu
④の図\ref{ph:波-7:fig-tsunagi}と同じ形。$Z_1$ は空気またはゼリー、$Z_2$ は体。
\phriyu
はね返る割合は $r^2$、入る割合は $4Z_1Z_2/(Z_1+Z_2)^2$（④）。インピーダンスの値だけで決まります。
\phtateru
\[r=\frac{Z_1-Z_2}{Z_1+Z_2},\qquad\text{はね返る割合}=r^2,\qquad\text{入る割合}=1-r^2\]
\phtoku
（1）\[r=\frac{413-\num{1630000}}{413+\num{1630000}}\approx-0.99949,\qquad r^2\approx0.9990\]
入るのは $1-0.9990=0.0010$、つまり0.1\,\% です。\\
（2）\[r=\frac{1.50-1.63}{1.50+1.63}=\frac{-0.13}{3.13}\approx-0.0415,\qquad r^2\approx0.0017\]
\phtashikame
単位：$Z$ の比なので単位はありません。\\
大きさ：空気と体の比は約4000倍で、ほとんど全部がはね返ります。ゼリーと体はほぼ同じ $Z$ なので、ほとんど全部が入ります。\\
極端な場合：$Z_1=Z_2$ なら $r=0$ で、まったくはね返りません。
\kotae （1）はね返る99.9\,\%、入る0.1\,\%。（2）はね返る0.17\,\%、入る99.8\,\%
\end{phreidai}
\begin{jibun}
体の中で、軟らかい組織（$Z=1.63\times10^6$）から、（1）骨（$7.8\times10^6$）、（2）肺の中の空気（413）へ音が進むとき、はね返るエネルギーの割合を求めましょう。超音波の検査で、骨の向こうや肺が写りにくいわけを考えましょう。
\bkjans{（1）$r\approx-0.65$、$r^2\approx43\,\%$。（2）$r^2\approx99.9\,\%$。どちらも、多くがはね返って奥へ届かない}
\end{jibun}

\label{ph:波-7:ex2}
\begin{phreidai}[音圧と音の強さ]
空気中を進む音の、音圧の実効値が\qty{1.0}{Pa}です。音の強さを求めましょう。また、音の大きさは何 dB ですか。音の大きさの dB は、音圧の実効値 $p_{\mathrm e}$ を\qty{20}{\micro\pascal}と比べた $20\log_{10}(p_{\mathrm e}/\qty{20}{\micro\pascal})$ とします（詳しくは波-9）。
\phbamen
平らな音の波（平面波）の強さを、音圧から求める場面です。
\phmushi
空気の温度は\qty{20}{\celsius}、吸収は無視します。
\phzu
$p_{\mathrm e}=\qty{1.0}{Pa}$、$Z=\rho c=\qty{413}{Pa.s/m}$。
\phriyu
音の強さは、圧力×速さの平均で、$I=p_{\mathrm e}^{\,2}/(\rho c)$（③）。
\phtateru
\[I=\frac{p_{\mathrm e}^{\,2}}{\rho c},\qquad L=20\log_{10}\frac{p_{\mathrm e}}{\qty{20}{\micro\pascal}}\]
\phtoku
\[I=\frac{1.0^2}{413}\approx\qty{2.4e-3}{W/m^2}\]
\[L=20\log_{10}\frac{1.0}{20\times10^{-6}}=20\log_{10}(5.0\times10^4)=20\times4.70\approx\qty{94}{dB}\]
\phtashikame
単位：Pa$^2$÷(Pa·s/m)＝Pa·m/s＝N/m$^2$×m/s＝W/m$^2$。\\
大きさ：\qty{94}{dB}は、すぐ近くを電車が通るくらいの、とても大きな音です。それでも、1 m$^2$ に届くエネルギーは1秒に約\qty{2}{mJ}で、とても小さいのです。空気の粒が動く速さ（実効値）は $1.0\div413\approx\qty{2.4}{mm/s}$ にすぎません。
\kotae 約 $\qty{2.4e-3}{W/m^2}$、約\qty{94}{dB}
\end{phreidai}
\begin{jibun}
（1）ふつうの会話の声（音圧の実効値\qty{0.020}{Pa}）の強さと dB を求めましょう。（2）水の中（$Z=1.48\times10^6$）で、音圧の実効値が\qty{1.0}{Pa}の音の強さを求め、例題2の何分の1か答えましょう。
\bkjans{（1）$0.020^2\div413$ より、約 $\qty{9.7e-7}{W/m^2}$。$20\log_{10}1000=\qty{60}{dB}$。（2）$1.0^2\div(1.48\times10^6)$ より、約 $\qty{6.8e-7}{W/m^2}$。約3600分の1}
\end{jibun}

\label{ph:波-7:ex3}
\begin{phreidai}[ロープをゆする手の仕事率]
線密度\qty{0.20}{kg/m}のロープを、張力\qty{80}{N}で張り、はしを振動数\qty{2.0}{Hz}、振幅\qty{0.050}{m}でゆすって、長いロープに正弦波を送ります。波の速さ、ロープのインピーダンス、手が1秒にする仕事（仕事率）を求めましょう。
\phbamen
手がロープにする仕事が、波のエネルギーとして運ばれていく場面です。
\phmushi
ロープは十分長く、先からはね返りは戻ってこないとします。ロープの中のこすれ（吸収）と重力は無視します。
\phzu
図\ref{ph:波-7:fig-rope}。$\mu=0.20$、$F_{\mathrm T}=80$、$f=2.0$、$A=0.050$。
\phriyu
手は、インピーダンス $Z$ のロープを速さ $u$ で動かすので、仕事率は $Zu^2$ の平均 $\tfrac12Z\omega^2A^2$（③）。
\phtateru
\[c=\sqrt{\frac{F_{\mathrm T}}{\mu}},\qquad Z=\mu c,\qquad P=\frac12Z\omega^2A^2\]
\phtoku
\[c=\sqrt{\frac{80}{0.20}}=\sqrt{400}=\qty{20}{m/s},\qquad Z=0.20\times20=\qty{4.0}{N.s/m}\]
\[\omega=2\pi\times2.0\approx\qty{12.6}{rad/s},\qquad P=\frac12\times4.0\times12.6^2\times0.050^2\approx\qty{0.79}{W}\]
\phtashikame
単位：N·s/m×(1/s)$^2$×m$^2$＝N·m/s＝W。\\
大きさ：①の式でも、$\overline E=\tfrac12\times0.20\times12.6^2\times0.050^2\approx\qty{0.039}{J/m}$、$P=\overline E\,c\approx\qty{0.79}{W}$ で合います。手の動きのいちばん速い所は $A\omega\approx\qty{0.63}{m/s}$ で、ゆっくりした動きですが、ずっと続けると腕が疲れます。\\
極端な場合：$A=0$ なら $P=0$。ゆすらなければ仕事はいりません。
\kotae $c=\qty{20}{m/s}$、$Z=\qty{4.0}{N.s/m}$、約\qty{0.79}{W}
\end{phreidai}
\begin{jibun}
例題3のロープを、（1）振動数\qty{1.0}{Hz}・振幅\qty{0.10}{m}、（2）振動数\qty{3.0}{Hz}・振幅\qty{0.050}{m}でゆするときの仕事率を求めましょう。
\bkjans{（1）約\qty{0.79}{W}（振動数が半分、振幅が2倍で同じ）。（2）$\tfrac12\times4.0\times(6\pi)^2\times0.050^2$ で、約\qty{1.8}{W}}
\end{jibun}

\label{ph:波-7:ex4}
\begin{phreidai}[線密度が4倍の弦につなぐ]
張力が同じで、線密度が4倍の弦をつなぎました。左から高さ\qty{3.0}{cm}の山を送ると、はね返る山と通りぬける山の高さはいくらですか。エネルギーの何分のいくつがはね返り、何分のいくつが通りぬけますか。
\phbamen
④の表計算と同じ、2本の弦のつなぎ目です。
\phmushi
結び目の質量、弦の中の吸収を無視します。
\phzu
図\ref{ph:波-7:fig-tsunagi}。線密度 $\mu_2=4\mu_1$、張力は同じ。
\phriyu
$Z=\sqrt{F_{\mathrm T}\mu}$ なので、張力が同じなら $Z$ は $\sqrt\mu$ に比例します。$Z$ の比がわかれば、$r$ と $\tau$ が決まります（④）。
\phtateru
\[\frac{Z_2}{Z_1}=\sqrt{\frac{\mu_2}{\mu_1}},\qquad r=\frac{Z_1-Z_2}{Z_1+Z_2},\qquad\tau=\frac{2Z_1}{Z_1+Z_2}\]
\phtoku
$Z_2/Z_1=\sqrt4=2$ なので、$Z_1=1$、$Z_2=2$ として、
\[r=\frac{1-2}{1+2}=-\frac13,\qquad\tau=\frac{2\times1}{1+2}=\frac23\]
山の高さは、はね返り $3.0\times(-\tfrac13)=\qty{-1.0}{cm}$、通りぬけ $3.0\times\tfrac23=\qty{2.0}{cm}$。エネルギーは、はね返り $r^2=\tfrac19$、通りぬけ $\tfrac{Z_2}{Z_1}\tau^2=2\times\tfrac49=\tfrac89$。
\phtashikame
エネルギーの和は $\tfrac19+\tfrac89=1$ で、保存しています。$1+r=\tau$（$1-\tfrac13=\tfrac23$）も成り立ちます。表計算（表\ref{ph:波-7:tab-sheet}・図\ref{ph:波-7:fig-sheet}）の結果とも合います。
\kotae はね返る山は反対側に\qty{1.0}{cm}、通りぬける山は\qty{2.0}{cm}。エネルギーは $\tfrac19$ がはね返り、$\tfrac89$ が通りぬける
\end{phreidai}
\begin{jibun}
張力が同じで、線密度が（1）$\tfrac14$ 倍、（2）9倍の弦につなぐとき、$r$、$\tau$ と、はね返るエネルギーの割合を求めましょう。
\bkjans{（1）$Z_2/Z_1=\tfrac12$、$r=+\tfrac13$、$\tau=\tfrac43$、はね返り $\tfrac19$（通りぬけ $\tfrac89$）。（2）$Z_2/Z_1=3$、$r=-\tfrac12$、$\tau=\tfrac12$、はね返り $\tfrac14$（通りぬけ $\tfrac34$）}
\end{jibun}

\label{ph:波-7:ex5}
\begin{phreidai}[空気から水へ入る音]
空気（$Z=413$）を伝わってきた音が、水面（$Z=1.48\times10^6$）に垂直に当たります。水の中へ入るエネルギーの割合を求め、dB で表しましょう。
\phbamen
やってみようの、お風呂の中で外の声が聞こえにくい場面です。
\phmushi
水面は平らで、波立っていないとします。音は水面に垂直に当たるとします。
\phzu
$Z_1=413$、$Z_2=1.48\times10^6$（Pa·s/m）。
\phriyu
入る割合は $4Z_1Z_2/(Z_1+Z_2)^2$（④）。割合を dB にするには $10\log_{10}$（3-12⑦）。
\phtateru
\[\text{入る割合}=\frac{4Z_1Z_2}{(Z_1+Z_2)^2},\qquad\text{dB}=10\log_{10}(\text{入る割合})\]
\phtoku
\[\frac{4\times413\times1.48\times10^6}{(413+1.48\times10^6)^2}\approx\frac{2.44\times10^9}{2.19\times10^{12}}\approx0.0011\]
\[10\log_{10}0.0011\approx\qty{-30}{dB}\]
\phtashikame
$Z_2\gg Z_1$ のときは、分母を $Z_2^{\,2}$ と近似して、
\[\text{入る割合}\approx\frac{4Z_1}{Z_2}=\frac{4\times413}{1.48\times10^6}\approx0.0011\]
と、すぐ見当がつきます。\\
大きさ：$-\qty{30}{dB}$ は、エネルギーで1000分の1です。外の会話の声が水面に当たっても、水の中へ入るエネルギーは、その1000分の1しかありません（水面のすぐ下の音圧は、⑤のせいかくに言うとのように、むしろ2倍近くになります。小さくなるのはエネルギーです）。\\
逆向き（水の中の音が空気へ出る）でも、同じ式なので、割合は同じ0.11\,\% です。
\kotae 約0.11\,\%（約 $\qty{-30}{dB}$）
\end{phreidai}
\begin{jibun}
水（$Z=1.48\times10^6$）から鋼（$Z=4.6\times10^7$）へ音が垂直に入るとき、はね返る割合と、入る割合を dB で求めましょう。
\bkjans{$r\approx-0.94$、はね返り $r^2\approx88\,\%$。入る割合 約12\,\%、約 $\qty{-9}{dB}$}
\end{jibun}

\label{ph:波-7:ex6}
\begin{phreidai}[超音波のプローブの整合層]
圧電セラミックス（$Z_1=3.0\times10^7$）から、ゼリー（$Z_2=1.5\times10^6$）へ、振動数\qty{5.0}{MHz}の超音波を出します。はね返りをなくす4分の1波長の整合層のインピーダンスと厚さを求めましょう。層の中の音速は\qty{3000}{m/s}とします。また、層がないときに、はね返るエネルギーの割合はいくらですか。
\phbamen
プローブの先にはる、薄い層を設計する場面です。
\phmushi
層の中の吸収を無視します。音は層に垂直に進み、層は一様な厚さとします。
\phzu
図\ref{ph:波-7:fig-seigou}。$Z_1=3.0\times10^7$、$Z_2=1.5\times10^6$、$c_{\mathrm m}=3000$、$f=5.0\times10^6$。
\phriyu
2つの面のはね返りが同じ高さで半波長ずれると打ち消し合う、という条件から $Z_{\mathrm m}=\sqrt{Z_1Z_2}$、厚さ $\lambda_{\mathrm m}/4$（⑧）。
\phtateru
\[Z_{\mathrm m}=\sqrt{Z_1Z_2},\qquad\lambda_{\mathrm m}=\frac{c_{\mathrm m}}{f},\qquad\text{厚さ}=\frac{\lambda_{\mathrm m}}{4}\]
\phtoku
\[Z_{\mathrm m}=\sqrt{3.0\times10^7\times1.5\times10^6}=\sqrt{4.5\times10^{13}}\approx6.7\times10^6\]
\[\lambda_{\mathrm m}=\frac{3000}{5.0\times10^6}=\qty{6.0e-4}{m}=\qty{0.60}{mm},\qquad\text{厚さ}=\qty{0.15}{mm}\]
層がないとき：$r=(30-1.5)/(30+1.5)\approx0.905$、$r^2\approx0.82$。
\phtashikame
$Z_{\mathrm m}$ は、$Z_1$ の約4.5分の1、$Z_2$ の約4.5倍で、2つのちょうど「かけ算のまん中」です。厚さ\qty{0.15}{mm}は、紙1枚くらいです。\\
極端な場合：振動数が変わると $\lambda_{\mathrm m}$ が変わり、ちょうど $\tfrac14$ でなくなるので、少しはね返ります（\qty{4.0}{MHz}では、通りぬけが約70\,\% に下がる）。
\kotae 約 $6.7\times10^6$ Pa·s/m、厚さ約\qty{0.15}{mm}。層がないと約82\,\% がはね返る
\end{phreidai}
\begin{jibun}
水（$Z=1.48\times10^6$）と鋼（$Z=4.6\times10^7$）のあいだに入れる4分の1波長の整合層のインピーダンスと、振動数\qty{2.0}{MHz}のときの厚さを求めましょう。層の中の音速は\qty{2500}{m/s}とします。
\bkjans{約 $8.3\times10^6$ Pa·s/m。$\lambda_{\mathrm m}=2500\div(2.0\times10^6)=\qty{1.25}{mm}$、厚さ約\qty{0.31}{mm}}
\end{jibun}

\label{ph:波-7:ex7}
\begin{phreidai}[50 Ω のケーブルに 75 Ω の機器をつなぐ]
特性インピーダンス\qty{50}{\ohm}の同軸ケーブルの先に、入力のインピーダンスが\qty{75}{\ohm}の機器をつなぎました。電圧の反射係数と、はね返る電力の割合を求めましょう。
\phbamen
電気の線を伝わる波が、インピーダンスの違う機器に入る場面です。
\phmushi
ケーブルの中での弱まりを無視します。機器のインピーダンスは、振動数によらず\qty{75}{\ohm}とします。
\phzu
$Z_{\text{線}}=\qty{50}{\ohm}$、$Z_{\text{先}}=\qty{75}{\ohm}$。
\phriyu
電圧は「押すもの」なので、反射係数は $(Z_{\text{先}}-Z_{\text{線}})/(Z_{\text{先}}+Z_{\text{線}})$（⑧）。はね返る電力の割合は、その2乗。
\phtateru
\[r_V=\frac{Z_{\text{先}}-Z_{\text{線}}}{Z_{\text{先}}+Z_{\text{線}}},\qquad\text{はね返る割合}=r_V^{\,2}\]
\phtoku
\[r_V=\frac{75-50}{75+50}=\frac{25}{125}=0.20,\qquad r_V^{\,2}=0.040\]
\phtashikame
はね返るのは4\,\% で、96\,\% は機器に入ります。テレビの受信なら問題にならない程度ですが、強い電波を出す無線機では、はね返った電力が無線機を傷めることがあるので、そろえます。\\
極端な場合：先をつながず開けておく（$Z_{\text{先}}\to\infty$）と $r_V=+1$、線の先をショートさせる（$Z_{\text{先}}=0$）と $r_V=-1$ で、全部はね返ります。
\kotae 電圧の反射係数0.20、はね返る電力4\,\%
\end{phreidai}
\begin{jibun}
（1）\qty{75}{\ohm}のケーブルに\qty{50}{\ohm}の機器をつなぐと、電圧の反射係数はいくらですか。（2）昔のテレビのアンテナ線（\qty{300}{\ohm}）を、\qty{75}{\ohm}の入り口に直接つなぐと、はね返る電力は何\,\% ですか。
\bkjans{（1）$-0.20$。（2）$r_V=-0.60$、はね返る電力 36\,\%}
\end{jibun}

\label{ph:波-7:ex8}
\begin{phreidai}[深い海の波の山と、エネルギーの速さ]
深い海で、波長\qty{10}{m}の波が進んでいます。山の進む速さ（位相速度）と、エネルギーの進む速さ（群速度）を求めましょう。分散関係は $\omega=\sqrt{gk}$、$g=\qty{9.8}{m/s^2}$ とします。
\phbamen
沖の波の、目に見える山の動きと、エネルギーの動きを比べる場面です。
\phmushi
波は小さく（波の高さは波長に比べて十分小さい）、深さは波長の半分より深く、風は無視します。
\phzu
図\ref{ph:波-7:fig-gun}。$\lambda=\qty{10}{m}$、$k=2\pi/\lambda$。
\phriyu
分散のある波では、山は $c=\omega/k$、エネルギーは群速度 $c_{\mathrm g}=\dd\omega/\dd k$ で進みます（⑥・波-5 ⑦）。
\phtateru
\[k=\frac{2\pi}{\lambda},\qquad c=\sqrt{\frac{g}{k}},\qquad c_{\mathrm g}=\frac{c}{2}\]
\phtoku
\[k=\frac{2\pi}{10}\approx\qty{0.628}{rad/m},\qquad c=\sqrt{\frac{9.8}{0.628}}\approx\qty{3.95}{m/s},\qquad c_{\mathrm g}\approx\qty{1.97}{m/s}\]
周期は $\lambda/c=10\div3.95\approx\qty{2.5}{s}$ です。
\phtashikame
単位：$\sqrt{(\text{m/s}^2)\div(\text{1/m})}=$ m/s。\\
大きさ：歩く速さ（約\qty{1.3}{m/s}）より少し速い山が、目に見えます。エネルギーはその半分の速さで届きます。\\
極端な場合：波長が長いほど $k$ が小さく、$c$ も $c_{\mathrm g}$ も大きくなります（波長\qty{100}{m}なら $c\approx\qty{12}{m/s}$）。
\kotae 位相速度 約\qty{3.9}{m/s}、群速度 約\qty{2.0}{m/s}
\end{phreidai}
\begin{jibun}
深い海で、周期\qty{10}{s}のうねりがあります。$\omega=\sqrt{gk}$ に $\omega=2\pi/T$、$k=2\pi/\lambda$ を入れると、$\lambda=gT^2/(2\pi)$、$c=gT/(2\pi)$ になります。うねりの波長、位相速度、群速度と、\qty{1000}{km}離れた台風の所から、エネルギーが届くまでの時間を求めましょう。
\bkjans{$\lambda\approx\qty{156}{m}$、$c\approx\qty{15.6}{m/s}$、$c_{\mathrm g}\approx\qty{7.8}{m/s}$。届くまで約 $\qty{1.28e5}{s}$、約36時間}
\end{jibun}

\label{ph:波-7:ex9}
\begin{phreidai}[超音波は、どこまで深く見えるか]
振動数\qty{5.0}{MHz}の超音波で、深さ\qty{10}{cm}の所を調べます。吸収を 1 cm・1 MHz あたり\qty{0.5}{dB}とすると、行って戻る音は何 dB 弱まりますか。強さで何分の1ですか。
\phbamen
体の中で、はね返って戻ってくる音の弱まりを見積もる場面です。
\phmushi
境目ではね返るときの割合と、音の広がりは考えず、吸収だけを数えます。
\phzu
深さ\qty{10}{cm}、往復の距離\qty{20}{cm}、$f=\qty{5.0}{MHz}$。
\phriyu
吸収の dB は、振動数と進む距離に比例してたし算できます（⑦）。
\phtateru
\[\text{弱まり〔dB〕}=0.5\times f\text{〔MHz〕}\times\text{往復の距離〔cm〕}\]
\phtoku
\[0.5\times5.0\times20=\qty{50}{dB},\qquad10^{-50/10}=10^{-5}\]
\phtashikame
\qty{50}{dB}は強さで10万分の1です。小さな戻りの音を、大きく増幅して絵にします。\\
極端な場合：深さ2倍の\qty{20}{cm}では\qty{100}{dB}になり、絵を作るのが難しくなります。
\kotae 約\qty{50}{dB}、強さで約10万分の1
\end{phreidai}
\begin{jibun}
（1）\qty{3.5}{MHz}で深さ\qty{15}{cm}を調べるとき、往復の弱まりは何 dB ですか。（2）\qty{10}{MHz}のプローブで、往復の弱まりを\qty{60}{dB}以下にしたいとき、調べられる深さは何 cm までですか。
\bkjans{（1）$0.5\times3.5\times30$ で、約\qty{53}{dB}。（2）往復\qty{12}{cm}なので、深さ\qty{6}{cm}まで}
\end{jibun}

\label{ph:波-7:ex10}
\begin{phreidai}[小さなスピーカーから、離れると]
小さなスピーカーが、1秒に\qty{0.010}{W}の音のエネルギーを、すべての向きに同じように出しています。\qty{1}{m}・\qty{2}{m}・\qty{10}{m}離れた所の音の強さを求め、強さ $10^{-12}$ W/m$^2$ を基準にした dB（$10\log_{10}(I/10^{-12})$）で表しましょう。
\phbamen
点の音源から、音が球面に広がっていく場面です。
\phmushi
床・壁ではね返る音と、空気の吸収を無視します。スピーカーは、距離に比べて十分小さいとします。
\phzu
図\ref{ph:波-7:fig-kyuu}。$P=\qty{0.010}{W}$、$R=1,\,2,\,10$ m。
\phriyu
吸収がなければ、どの球面を通るエネルギーも同じなので、$I=P/(4\pi R^2)$（②）。
\phtateru
\[I=\frac{P}{4\pi R^2},\qquad L=10\log_{10}\frac{I}{10^{-12}}\]
\phtoku
\[R=\qty{1}{m}:\ I=\frac{0.010}{4\pi\times1^2}\approx\qty{8.0e-4}{W/m^2},\quad L=10\log_{10}(8.0\times10^8)\approx\qty{89}{dB}\]
$R=\qty{2}{m}$：$I$ は $\tfrac14$ の $\qty{2.0e-4}{W/m^2}$、$L\approx89-6=\qty{83}{dB}$。\\
$R=\qty{10}{m}$：$I$ は $\tfrac1{100}$ の $\qty{8.0e-6}{W/m^2}$、$L\approx89-20=\qty{69}{dB}$。
\phtashikame
距離2倍で約\qty{6}{dB}、10倍で\qty{20}{dB}下がる、という②の決まりに合います。\\
大きさ：\qty{0.010}{W}の音は、かなり大きな音です（スピーカーの電気の入力は、その何十倍もある）。
\kotae \qty{1}{m}：約 $\qty{8.0e-4}{W/m^2}$（\qty{89}{dB}）、\qty{2}{m}：約 $\qty{2.0e-4}{W/m^2}$（\qty{83}{dB}）、\qty{10}{m}：約 $\qty{8.0e-6}{W/m^2}$（\qty{69}{dB}）
\end{phreidai}
\begin{jibun}
点の音源から\qty{5}{m}の所で\qty{70}{dB}でした。\qty{20}{m}と\qty{50}{m}の所では何 dB ですか（吸収とはね返りは無視）。
\bkjans{\qty{20}{m}は距離4倍で、約\qty{58}{dB}。\qty{50}{m}は10倍で、\qty{50}{dB}}
\end{jibun}

\label{ph:波-7:ex11}
\begin{phreidai}[はね返らない端を作る]
例題3のロープ（$Z=\qty{4.0}{N.s/m}$）の先を、ダンパー（速さ $u$ に比例する力 $bu$ で止める筒）につなぎます。はね返りがなくなる $b$ を求めましょう。また、$b=\qty{6.0}{N.s/m}$ のダンパーでは、はね返るエネルギーの割合はいくらですか。
\phbamen
実験室で、ロープの先からのはね返りを止めたい場面です（長いロープの代わり）。
\phmushi
ダンパーの先の金具の質量、ダンパーの中のばねの働きを無視します。
\phzu
ロープの先 $x=0$ にダンパー。$Z_1=4.0$、$Z_{\text{端}}=b$。
\phriyu
ダンパーは、はしの速さに比例する力で押し返すので、インピーダンス $b$ の「弦2」と同じ働きをします（⑤）。
\phtateru
\[r=\frac{Z_1-b}{Z_1+b}\]
\phtoku
$r=0$ になるのは $b=Z_1=\qty{4.0}{N.s/m}$ のとき。$b=6.0$ では、
\[r=\frac{4.0-6.0}{4.0+6.0}=-0.20,\qquad r^2=0.040\]
\phtashikame
$b=Z_1$ のダンパーは、ロープから見ると「どこまでも続くロープ」と見分けがつきません。来たエネルギーは全部ダンパーの中で熱になります。$b$ が少しずれても、はね返りは4\,\% と小さく済みます。\\
極端な場合：$b\to\infty$ で固定端、$b=0$ で自由端にもどります。
\kotae $b=\qty{4.0}{N.s/m}$。$b=\qty{6.0}{N.s/m}$ では約4\,\%
\end{phreidai}
\begin{jibun}
（1）例題3のロープの先に、$b=\qty{2.0}{N.s/m}$ のダンパーをつけると、$r$ とはね返る割合はいくらですか。（2）ダンパーの代わりに、ばね定数\qty{8.0}{N/m}のばねをつけ、\qty{1.0}{Hz}でゆすると、$|r|$ と $r$ の角はいくらですか（⑤）。
\bkjans{（1）$r=(4.0-2.0)\div(4.0+2.0)=+\tfrac13$、約11\,\%。（2）$|r|=1$。角は $2\arctan(8.0\div25.1)$ で、約35°}
\end{jibun}

\label{ph:波-7:ex12}
\begin{phreidai}[ガラスの表面で、はね返る光]
空気から、屈折率1.5のガラスへ、光が垂直に当たります。1つの面ではね返る光の割合を求めましょう。また、窓ガラス1枚（表と裏の2面）では、合わせて約何\,\% がはね返りますか。
\phbamen
窓ガラスやメガネのレンズの表面での、光のはね返りです。
\phmushi
ガラスの中での吸収を無視します。光はガラスの面に垂直に当たるとします。
\phzu
$n_1=1.00$、$n_2=1.5$。
\phriyu
光のインピーダンスは屈折率に反比例するので、はね返る割合は、屈折率だけで書いた⑨の式で求められます。
\phtateru
\[r^2=\Bigl(\frac{n_1-n_2}{n_1+n_2}\Bigr)^2\]
\phtoku
\[r^2=\Bigl(\frac{1.00-1.5}{1.00+1.5}\Bigr)^2=\Bigl(\frac{-0.5}{2.5}\Bigr)^2=0.040\]
裏の面（ガラス→空気）でも、同じ式で4\,\% です。表で96\,\% が入り、その4\,\% が裏ではね返るので、2面で約 $4+0.96\times4\approx7.8$\,\%（ガラスの中を何回も行き来する分まで入れると約7.7\,\%）。
\phtashikame
4\,\% は、夜の窓に部屋の中が写るくらいの明るさです（⑨の身近な例）。\\
極端な場合：$n_2=n_1$ ならはね返りは0です。水の中に入れたガラスは、$n$ の差が小さいので、ほとんど見えなくなります。
\kotae 1面で4\,\%、窓ガラス1枚で約8\,\%
\end{phreidai}
\begin{jibun}
（1）ダイヤモンド（$n=2.42$）の表面で、垂直に当たった光がはね返る割合を求めましょう。（2）水（$n=1.333$）の中に入れたガラス（$n=1.5$）の表面ではどうですか。
\bkjans{（1）約17\,\%。（2）約0.35\,\%（ほとんどはね返らないので、見えにくい）}
\end{jibun}

%% ======================== まちがえやすいところ ========================
\Needspace{11\baselineskip}\bksection{tsumazuki}{まちがえやすいところをチェック}
\begin{tsumazuki}{振幅が2倍なら、何 dB 上がる？}
\bkng{振幅が2倍なので、$10\log_{10}2\approx\qty{3}{dB}$ 上がる。}
\bkok{強さは振幅の2乗に比例するので4倍。$10\log_{10}4\approx\qty{6}{dB}$ 上がる（$20\log_{10}2$ と同じ）。}
dB の $10\log_{10}$ は、強さ（エネルギー）の比に使います。振幅や音圧の比なら $20\log_{10}$ です（3-12⑦、①）。
\end{tsumazuki}

\begin{tsumazuki}{反射係数と、はね返る割合}
\bkng{$r=-\tfrac13$ なので、エネルギーの3分の1がはね返る。}
\bkok{はね返るエネルギーの割合は $r^2=\tfrac19$。符号は、山が反対側になることを表す。}
$r$ は高さの比、$r^2$ はエネルギーの比です（④）。
\end{tsumazuki}

\begin{tsumazuki}{通りぬけた波のほうが高い？}
\bkng{軽い弦へ通りぬけた山（$\tau=\tfrac43$）のほうが、来た山より高いのはおかしい。エネルギーが増えている。}
\bkok{山は高くなるが、$Z$ が小さい弦なので、エネルギーは $\tfrac{Z_2}{Z_1}\tau^2=\tfrac12\times\tfrac{16}{9}=\tfrac89$ に減っている。}
エネルギーは高さだけでなく、$Z$ にも比例します（③）。
\end{tsumazuki}

\begin{tsumazuki}{速い物ほど、インピーダンスが大きい？}
\bkng{波が速く伝わる物ほど、$Z$ が大きい。}
\bkok{$Z=\rho c$（弦では $\mu c$）なので、密度も効く。弦の線密度を4倍にすると、$c$ は半分になるが、$Z$ は2倍になる。}
$Z$ と $c$ は別の量です。はね返りを決めるのは $Z$ の比です（③）。
\end{tsumazuki}

\begin{tsumazuki}{見えている山の速さで、エネルギーが届く？}
\bkng{沖の波の山が\qty{4}{m/s}で進んでいるので、エネルギーも\qty{4}{m/s}で岸に届く。}
\bkok{深い海の波では、エネルギーは群速度で進み、山の速さの半分の約\qty{2}{m/s}。}
分散のある波では、山の速さとエネルギーの速さが違います（⑥）。
\end{tsumazuki}

%% 修3-4（P4-14）：工学ではこう使うの囲みは分けられず、次のページへ送られて下の空きが半分になったので、コラムをここへ移した
\begin{column}{はじめの超音波の検査は、水の中で}
超音波で体の中を見る研究は、1942年、オーストリアのドゥシックが、頭に超音波を通して脳の病気を見つけようとした実験から始まり、1950年代に本格的に広がりました。アメリカのハウリーらは、はじめ、調べる人を水の入った大きなおけに入れ、水の中から超音波を当てる装置を作りました。空気の層をなくすために、体を水に沈めたのです。イギリスのドナルドらは、工場で金属の傷を探す超音波探傷の装置を借りて、おなかの中を調べ、1958年に論文を出しました。その後、プローブを直接体に当て、すきまをゼリーで満たすやり方が広まり、今の検査の形になりました。体を水に沈めるのも、ゼリーをぬるのも、目的は同じ「空気を追い出して、$Z$ をそろえる」ことです。
\end{column}

%% ======================== 工学ではこう使う ========================
\Needspace{18\baselineskip}\bksection{phgear}{工学ではこう使う}
\begin{phkougaku}[はね返りを使う・はね返りを減らす]
\textbf{超音波の検査と探傷}：超音波の検査では、体の中の組織の境目で、わずかにはね返った音（$r^2$ は0.1〜1\,\% くらい）を集めて、絵を作ります。音が戻るまでの時間から境目の深さを、戻った音の強さから境目の $Z$ の違いを読みとります（波-10）。鉄の橋やレールの中の割れ目（傷）を探す\textbf{超音波探傷}では、鋼（$4.6\times10^7$）と割れ目の中の空気（413）の $Z$ がとても違うので、割れ目でほとんど全部がはね返り、小さな割れ目でもはっきり見つかります。

\textbf{魚群探知機}：船の底から海の中へ音を出し、魚の群れや海の底からはね返る音を聞きます。魚の体の中の浮き袋（空気）は、水と $Z$ が大きく違うので、よくはね返ります。

\textbf{レンズとケーブル}：メガネやカメラのレンズの反射防止膜、テレビのアンテナ線と機器のインピーダンスの統一（\qty{75}{\ohm}）は、どれも「$Z$ をそろえて、はね返りを減らす」整合です。

\begin{phketa}[インピーダンスと反射の数の目安]
\phketaline{音響インピーダンス（Pa·s/m）}{空気 約400、水・体 約 $1.5\times10^6$、鋼 約 $5\times10^7$}
\phketaline{はね返る割合}{空気→水・体 99.9\,\%、水→鋼 約90\,\%、空気→ガラス（光）4\,\%}
\phketaline{体の中の吸収}{約 \mbox{0.5\,dB/(cm·MHz)}。5 MHz・深さ\qty{10}{cm}の往復で約\qty{50}{dB}}
\phboxhead{答えが桁で合っているか}
はね返る割合は、0から1のあいだです。$r$ の大きさが1をこえたり、はね返る割合と通りぬける割合の和が1にならなかったりしたら、$Z_1$ と $Z_2$ を入れまちがえたか、$r$ と $r^2$ を取り違えています。$Z$ の比が10倍をこえると、半分以上がはね返る、と覚えておくと便利です。
\end{phketa}
\end{phkougaku}
\phstation{医用超音波,非破壊検査,音響工学,高周波回路}{音響インピーダンス、反射係数、透過率、整合層、減衰定数（dB/cm）、群速度、特性インピーダンス}

%% ======================== 章末問題 ========================
\bkendproblems
\bkproblemgroup{基本}{3}
\bkproblem 線密度\qty{0.30}{kg/m}のロープを張力\qty{120}{N}で張りました。（1）波の速さとインピーダンスを求めましょう。（2）はしを振動数\qty{1.5}{Hz}、振幅\qty{0.080}{m}でゆすって長いロープに正弦波を送るとき、手の仕事率を求めましょう。
\bkpans{（1）$c=\qty{20}{m/s}$、$Z=\qty{6.0}{N.s/m}$。（2）$\tfrac12\times6.0\times(3\pi)^2\times0.080^2$ で、約\qty{1.7}{W}}
\bkproblem 点の音源から\qty{1}{m}の所で\qty{90}{dB}でした。（1）\qty{3}{m}の所では何 dB ですか。（2）\qty{70}{dB}になるのは何 m の所ですか（吸収とはね返りは無視）。
\bkpans{（1）$90-20\log_{10}3$ で、約\qty{80}{dB}。（2）\qty{20}{dB}下がるのは距離10倍なので、\qty{10}{m}}
\bkproblem 空気中の音の、音圧の実効値が\qty{2.0}{Pa}です。音の強さと、音の大きさ（\qty{20}{\micro\pascal}を基準にした dB）を求めましょう。
\bkpans{$I=2.0^2\div413$ で、約 $\qty{9.7e-3}{W/m^2}$。$20\log_{10}(1.0\times10^5)=\qty{100}{dB}$}
\bkproblem （1）空気から水へ入る音のエネルギーの割合は約0.11\,\% でした（例題5）。逆に、水の中の音が水面から空気へ出ていくとき、出ていく割合はいくらですか。（2）体（$1.63\times10^6$）と水（$1.48\times10^6$）の境目で、はね返るエネルギーの割合はいくらですか。
\bkpans{（1）同じ式なので約0.11\,\%。（2）$r\approx-0.048$、$r^2\approx0.23\,\%$}
\bkproblem 張力が同じで、線密度が $\tfrac19$ 倍の弦につなぎました。高さ\qty{4.0}{cm}の山を送ると、はね返る山と通りぬける山の高さはいくらですか。はね返るエネルギーの割合も求めましょう。
\bkpans{$Z_2/Z_1=\tfrac13$、$r=+\tfrac12$、$\tau=\tfrac32$。はね返る山は同じ側に\qty{2.0}{cm}、通りぬける山は\qty{6.0}{cm}。はね返る割合 $\tfrac14$}
\bkproblemgroup{チャレンジ}{5}
\bkproblem 例題3のロープ（$Z_1=\qty{4.0}{N.s/m}$）の先に、ばね定数\qty{8.0}{N/m}のばねをつけます。はね返る波の角（$r$ の角）がちょうど90°になる振動数を求めましょう。
\bkpans{角が90°になるのは $k_{\mathrm s}/(\omega Z_1)=1$ のとき。$\omega=8.0\div4.0=\qty{2.0}{rad/s}$ なので、約\qty{0.32}{Hz}}
\bkproblem 深い海で、\qty{1000}{km}離れた台風の所から、周期\qty{15}{s}のうねりと周期\qty{10}{s}のうねりが同時に出ました。それぞれのエネルギーが届くまでの時間と、その差を求めましょう（群速度 $c_{\mathrm g}=gT/(4\pi)$）。
\bkpans{周期\qty{15}{s}：$c_{\mathrm g}\approx\qty{11.7}{m/s}$、約24時間。周期\qty{10}{s}：約\qty{7.8}{m/s}、約36時間（35.6）。差は約12時間}
\bkproblem 超音波の装置が、往復で\qty{80}{dB}の弱まりまで絵にできるとします。吸収を、1 cm・1 MHz あたり\qty{0.5}{dB}として、\qty{3.5}{MHz}と\qty{7.0}{MHz}のプローブで調べられる深さを求めましょう。振動数を2倍にすると、何がよくなり、何が悪くなりますか。
\bkpans{\qty{3.5}{MHz}：$80\div(0.5\times3.5\times2)\approx22.9$ で、約\qty{22}{cm}まで（届く深さの上限なので切り捨てる）。\qty{7.0}{MHz}：約\qty{11}{cm}。2倍にすると、波長が半分で細かく見えるが、届く深さは半分になる}
\bkproblem 振動数\qty{5.0}{MHz}の超音波を、ゼリーを通して体に入れ、深さ\qty{4.0}{cm}の骨ではね返った音を、同じプローブで受けとります。ゼリー（$Z=1.5\times10^6$）と体（$1.63\times10^6$）の境目を行きと帰りに通りぬける割合、体と骨（$7.8\times10^6$）の境目ではね返る割合、体の中の吸収（1 cm・1 MHz あたり\qty{0.5}{dB}）を合わせると、戻ってくる音の強さは、出した音より何 dB 小さいですか。強さで約何分の1ですか。3つのうち、どれがいちばん効きますか（例題1・例題9 と、じぶんでとこう1 を組み合わせる）。
\bkpans{通りぬける割合は $1-0.0017$ で、2回で約 $-0.015$\,dB。骨ではね返る割合は $r=(1.63-7.8)$ $\div(1.63+7.8)$ $\approx-0.65$、$r^2\approx0.43$ で約 $-3.7$\,dB。吸収は $0.5\times5.0\times8.0=\qty{20}{dB}$。合わせて約\qty{24}{dB}小さい（強さで約230分の1）。吸収がいちばん効く}
\bkproblem メガネのレンズ（屈折率1.50）の表面に、空気の中で波長\qty{550}{nm}の緑の光がはね返らないようにする4分の1波長の膜をつけます（⑧⑨）。光では $Z\propto1/n$ なので、$Z_{\mathrm m}=\sqrt{Z_1Z_2}$ は $n_{\mathrm m}=\sqrt{n_1n_2}$ になります。膜の中の波長は $\lambda/n_{\mathrm m}$ です。（1）理想の膜の屈折率と厚さを求めましょう。（2）フッ化マグネシウム（$n=1.38$）の膜なら、厚さはいくらにしますか。
\bkpans{（1）$n_{\mathrm m}=\sqrt{1.00\times1.50}\approx1.22$。厚さ $550\div(4\times1.22)\approx\qty{112}{nm}$。（2）$550\div(4\times1.38)\approx\qty{100}{nm}$}
\bkproblem ④の表計算のつなぎ目の式を導きましょう。つなぎ目の点に、左の弦の半マス分の質量 $\tfrac12\mu_1\Delta x_1$ と右の弦の半マス分の質量 $\tfrac12\mu_2\Delta x_2$ をもたせ、両どなりの点から張力で引かれる運動方程式を、$\Delta t$ ごとの式（波-5 ⑧と同じ形）にします。$\Delta x_1=c_1\Delta t$、$\Delta x_2=c_2\Delta t$ とし、$F_{\mathrm T}/c=Z$ を使います。
\bkpans{質量は $\tfrac12(\mu_1c_1+\mu_2c_2)\Delta t$、つまり $\tfrac12(Z_1+Z_2)\Delta t$。たての力は、$F_{\mathrm T}/\Delta x_2=Z_2/\Delta t$、$F_{\mathrm T}/\Delta x_1=Z_1/\Delta t$ より、$\{Z_2(y_{\text{右}}-y)+Z_1(y_{\text{左}}-y)\}/\Delta t$。これを $y^{\text{次}}-2y+y^{\text{前}}=\Delta t^2\times$（力÷質量）に入れて整理すると、$y^{\text{次}}=-y^{\text{前}}$ $+\{2Z_2y_{\text{右}}+2Z_1y_{\text{左}}\}/(Z_1+Z_2)$}

\begin{furikaeri}
\begin{checklist}
\item 波が運ぶエネルギーと強さを、振幅・振動数・距離から求められる。\bkfurref{ph:波-7:sec1}
\item インピーダンスの比から、境目ではね返る割合と通りぬける割合を求められる。\bkfurref{ph:波-7:sec4}
\item 表計算で、2本の弦をつないだときの山のはね返りと通りぬけを追える。\bkfurref{ph:波-7:sec4}
\item 波の式と境界条件から反射係数を導き、固定端・自由端・はね返らない端を1つの式で説明できる。\bkfurref{ph:波-7:sec5}
\item 吸収による弱まり、整合層、光の反射率を、インピーダンスと dB で見積もれる。\bkfurref{ph:波-7:sec7}
\end{checklist}
\end{furikaeri}
\bkfukidashi{mina}{ゼリーは、空気を追い出して、インピーダンスをそろえるためだったんだね。耳の中の小さな骨も、メガネの膜も、同じ考えでできているなんて！}
